\documentclass[10pt]{article}
\usepackage[T1,T2A]{fontenc}
\usepackage[utf8]{inputenc}
\usepackage[english,ukrainian]{babel}
\usepackage{geometry}
\usepackage{longtable}
\usepackage{array}
\usepackage{multirow}
\usepackage{hyperref}
\usepackage{mathptmx}
\renewcommand{\familydefault}{\rmdefault}

\geometry{a4paper,left=2cm,right=2cm,top=2cm,bottom=2cm}

\begin{document}

\begin{flushright}
ЗАТВЕРДЖЕНО\\
Наказом Відповідного Органу\\
від \_\_ \_\_\_\_\_\_\_\_\_ 2026 р. № \_\_\_\\
(в редакції наказу Відповідного Органу\\
від \_\_ \_\_\_\_\_\_\_\_\_ 2026 р. № \_\_)
\end{flushright}

\vspace{2cm}

\begin{center}
\textbf{\Large ІНФОРМАЦІЙНО-КОМУНІКАЦІЙНА СИСТЕМА}\\
\textbf{\Large «L1»}\\
\vspace{1cm}
\textbf{\Large ЦІЛЬОВИЙ ПРОФІЛЬ БЕЗПЕКИ}\\
\vspace{0.5cm}
\large UA.43220275.СЗІ.ЦПБ-01
\end{center}

\newpage

\begin{longtable}{|c|>{\raggedright\arraybackslash}p{2.5cm}|>{\raggedright\arraybackslash}p{4.5cm}|c|>{\raggedright\arraybackslash}p{4.5cm}|}
\hline
\multirow{2}{*}{\textbf{№}} & \multirow{2}{*}{\textbf{Вимога з безпеки інформації}} & \multirow{2}{*}{\textbf{Вимога ГПБ}} & \multicolumn{2}{c|}{\textbf{ЦПБ}} \\
\cline{4-5}
& & & \textbf{Захід захисту} & \textbf{Налаштований зміст заходу захисту}\\
\hline
\endfirsthead

\hline
\multirow{2}{*}{\textbf{№}} & \multirow{2}{*}{\textbf{Вимога з безпеки інформації}} & \multirow{2}{*}{\textbf{Вимога ГПБ}} & \multicolumn{2}{c|}{\textbf{ЦПБ}} \\
\cline{4-5}
& & & \textbf{Захід захисту} & \textbf{Налаштований зміст заходу захисту}\\
\hline
\endhead

\hline
\endfoot

\hline
\endlastfoot

\multicolumn{5}{|>{\raggedright\arraybackslash}p{14.5cm}|}{\textbf{1. УПРАВЛІННЯ ДОСТУПОМ (AC) (AC)}} \\ \hline
1 & id-spe-ac-1 &  & AC-1 & Визначено персонал або ролі, на які поширюється політика контролю доступу = \textbf{adminsecurity\_\allowbreak{}officer}\newline Визначено персонал або ролі, на які поширюються процедури контролю доступу = \textbf{adminsecurity\_\allowbreak{}officer}\newline Визначено посадову особу, яка керуватиме політикою та процедурами контролю доступу = \textbf{adminsecurity\_\allowbreak{}officer}\newline Визначено частоту, з якою переглядається та оновлюється поточна політика контролю доступу = \textbf{default\_\allowbreak{}deny\_\allowbreak{}ruleabac\_\allowbreak{}rule\_\allowbreak{}1}\newline Визначено події, які вимагають перегляду та оновлення поточної політики контролю доступу = \textbf{default\_\allowbreak{}deny\_\allowbreak{}ruleabac\_\allowbreak{}rule\_\allowbreak{}1}\newline Визначено частоту, з якою переглядаються та оновлюються поточні процедури контролю доступу = \textbf{30} \\ \hline
2 & УПРАВЛІННЯ ОБЛІКОВИМИ ЗАПИСАМИ (AC-2) &  & AC-2 & Визначено передумови та критерії членства в групах і ролях; визначено атрибути (за необхідності) для кожного облікового запису; визначено персонал або ролі, необхідні для затвердження запитів на створення облікових записів; визначено політику, процедури, передумови та критерії створення, активації, зміни, деактивації та видалення облікових записів; визначено персонал або ролі, які мають бути повідомлені; визначено період часу, протягом якого адміністратори облікових записів повинні бути повідомлені про те, що облікові записи більше не потрібні; визначено термін, протягом якого необхідно повідомляти адміністраторів облікових записів про звільнення або переведення користувачів; визначено період часу, протягом якого необхідно повідомляти адміністраторів облікових записів про зміни у використанні системи або необхідність знати про зміни для окремої особи; визначено атрибути, необхідні для авторизації доступу до системи (за потреби); AC-02\_\allowbreak{}ODP[10] AC-02a.[01] AC-02a.[02] AC-02b AC-02с AC-02d.01 AC-02d.02 AC-02d.03[01] AC-02d.03[02] AC-02e AC-02f.[01] AC-02f.[02] AC-02f.[03] AC-02f.[04] AC-02f.[05] AC-02g AC-02h.01 AC-02h.02 AC-02h.03 AC-02i.01 AC-02i.02 AC-02i.03 визначено періодичність перегляду облікових записів; визначено та задокументовано типи облікових записів, дозволених для використання в системі; визначено та задокументовано типи облікових записів, які заборонено використовувати в системі; призначені менеджери облікових записів; необхідні умови та критерії для членства в групах та ролях; визначено авторизованих користувачів системи; вказано приналежність до групи або ролі; для кожного облікового запису вказуються повноваження доступу (тобто привілеї); атрибути (за необхідності) вказуються для кожного облікового запису; для запитів на створення облікових записів потрібні схвалення від персоналу або ролей; облікові записи створюються відповідно до політики, процедур, передумов та критеріїв; облікові записи активуються відповідно до політики, процедур, передумов та критеріїв; облікові записи змінюються відповідно до політики, процедур, передумов та критеріїв; облікові записи деактивуються відповідно до політики, процедур, передумов та критеріїв; облікові записи видаляються відповідно до політики, процедур, передумов та критеріїв; контролюється використання облікових записів; адміністратори облікових записів та персонал або ролі отримують повідомлення протягом періоду часу, коли облікові записи більше не потрібні; адміністратори облікових записів та персонал або ролі отримують повідомлення протягом періоду часу, коли користувачі звільнені чи переведені; адміністратори облікових записів та персонал або ролі отримують повідомлення протягом періоду часу, коли використовуються індивідуальні системи або наявні зміни, які потребують нових знань. доступ до системи здійснюється на підставі дійсної авторизації доступу; доступ до системи авторизується на основі передбачуваного використання системи; доступ до системи авторизовано на основі атрибутів (за необхідності); AC-02j AC-02k.[01] AC-02k.[02] AC-02l.[01] AC-02l.[02] облікові записи переглядаються на відповідність вимогам управління обліковими записами частота; створено процес повторного випуску облікових даних спільного доступу або групових облікових записів (якщо вони розгорнуті), коли користувачів вилучено з групи; впроваджено процес повторного випуску облікових даних спільного доступу або групових облікових записів (якщо вони розгорнуті), коли користувачів вилучено з групи; процеси управління обліковими записами узгоджуються з процесами звільнення персоналу; процеси управління обліковими записами узгоджуються з процесами переводу персоналу = \textbf{24} \\ \hline
3 & id-spe-ac-2-3 &  & AC-2(3) & Облікові записи деактивуються протягом часового періоду, коли облікові записи більше не пов'язані з користувачем або фізичною особою = \textbf{adminsecurity\_\allowbreak{}officer}\newline Облікові записи відключено протягом часового періоду, коли облікові записи порушують політику організації; облікові записи деактивуються протягом , якщо вони були неактивними протягом  = \textbf{default\_\allowbreak{}deny\_\allowbreak{}ruleabac\_\allowbreak{}rule\_\allowbreak{}1}\newline Визначено період часу, протягом якого необхідно деактивувати облікові записи = \textbf{30}\newline Визначено період часу неактивності, після закінчення якого облікові записи будуть деактивовані; AC-02(03)(a) облікові записи деактивуються протягом часового періоду, коли термін дії облікових записів минув = \textbf{loginlogoutfailed\_\allowbreak{}attempt} \\ \hline
4 & id-spe-ac-2-5 &  & AC-2(5) & Користувачі повинні виходити з системи, коли період очікуваної бездіяльності або опис часу, коли потрібно вийти з системи = \textbf{30}\newline Визначено часовий період очікуваної бездіяльності або опис, коли потрібно вийти з системи = \textbf{30} \\ \hline
5 & id-spe-ac-2-13 &  & AC-2(13) & Облікові записи користувачів деактивуються протягом періоду часу з моменту виявлення значних ризиків = \textbf{30}\newline Визначено період часу, протягом якого необхідно деактивувати облікові записи фізичних осіб, які становлять значний ризик = \textbf{30} \\ \hline
6 & ЗАБЕЗПЕЧЕННЯ ДОСТУПУ (AC-3) &  & AC-3 & Затверджені повноваження на логічний доступ до інформації та ресурсів системи виконуються відповідно до чинних політик(правил) управління доступом = \textbf{default\_\allowbreak{}deny\_\allowbreak{}ruleabac\_\allowbreak{}rule\_\allowbreak{}1} \\ \hline
7 & УПРАВЛІННЯ ІНФОРМАЦІЙНИМИ ПОТОКАМИ (AC-4) &  & AC-4 & Затверджені повноваження застосовуються для контролю потоку інформації всередині системи та між підключеними системами на основі політики управління інформаційними потоками = \textbf{default\_\allowbreak{}deny\_\allowbreak{}ruleabac\_\allowbreak{}rule\_\allowbreak{}1}\newline Визначено політики управління інформаційними потоками всередині системи та між підключеними системами = \textbf{default\_\allowbreak{}deny\_\allowbreak{}ruleabac\_\allowbreak{}rule\_\allowbreak{}1} \\ \hline
8 & РОЗМЕЖУВАННЯ ОБОВ'ЯЗКІВ (AC-5) &  & AC-5 &  \\ \hline
9 & МІНІМІЗАЦІЯ ПОВНОВАЖЕНЬ (AC-6) &  & AC-6 &  \\ \hline
10 & id-spe-ac-6-1 &  & AC-6(1) & Визначені особи або ролі з авторизованим доступом до функцій безпеки та інформації, що має відношення до безпеки = \textbf{adminsecurity\_\allowbreak{}officer} \\ \hline
11 & id-spe-ac-6-2 &  & AC-6(2) & Користувачі облікових записів (або ролей) системи з доступом до функцій безпеки або інформації, що стосується безпеки, повинні використовувати непривілейовані облікові записи або ролі під час доступу до незахищених функцій = \textbf{adminsecurity\_\allowbreak{}officer} \\ \hline
12 & id-spe-ac-6-5 &  & AC-6(5) & Привілейовані облікові записи в системі обмежено персонал або ролі = \textbf{adminsecurity\_\allowbreak{}officer}\newline Визначено персонал або ролі, яким мають бути обмежені привілейовані облікові записи в системі = \textbf{adminsecurity\_\allowbreak{}officer} \\ \hline
13 & id-spe-ac-6-7 &  & AC-6(7) & Повноваження, призначені ролям і класам, переглядаються з частотою для перевірки необхідності таких повноважень = \textbf{30}\newline Визначено частоту перегляду повноважень, призначених ролям або класам користувачів = \textbf{30}\newline Визначено ролі або класи користувачів, яким призначено повноваження  = \textbf{adminsecurity\_\allowbreak{}officer} \\ \hline
14 & id-spe-ac-6-9 &  & AC-6(9) &  \\ \hline
15 & id-spe-ac-6-10 &  & AC-6(10) &  \\ \hline
16 & НЕВДАЛІ СПРОБИ ВХОДУ В СИСТЕМУ (AC-7) &  & AC-7 & Застосовано обмеження на кількість послідовних неуспішних спроб входу користувача протягом періоду часу = \textbf{30}\newline Визначається кількість послідовних неуспішних спроб входу користувача, дозволених протягом певного періоду часу = \textbf{30}\newline Визначено період часу, яким обмежується кількість послідовних неуспішних спроб входу користувача = \textbf{30}\newline Вибрано одне або декілька з наступних ЗНАЧЕНЬ ПАРАМЕТРІВ: \{заблокувати обліковий запис або вузол на період часу; заблокувати обліковий запис або вузол до зняття адміністратором; затримати наступний запит на вхід за алгоритмом затримки; повідомити системного адміністратора; виконати іншу дію\} = \textbf{AES-256-GCM}\newline Період часу, на який буде заблоковано обліковий запис або вузол (якщо вибрано) = \textbf{30}\newline Визначено алгоритм затримки наступного запиту на вхід (якщо вибрано) = \textbf{AES-256-GCM}\newline Інша дія, яка буде виконана після перевищення максимальної кількості невдалих спроб (якщо вибрано) = \textbf{3} \\ \hline
17 & id-spe-ac-8 &  & AC-8 & Визначається сповіщення або банер про використання системи, який система буде показувати користувачам перед наданням доступу до системи; визначено умови використання системи, які система відображатиме перед наданням подальшого доступу; повідомлення про використання системи відображається користувачам перед наданням доступу до системи, яке містить повідомлення про конфіденційність і безпеку, що відповідають чинним законам, нормативним документам, наказам, директивам, політикам, правилам, стандартам і керівним принципам; у повідомленні про використання системи зазначено, що користувачі отримують доступ до урядової системи; у повідомленні про використання системи зазначено, що використання системи може контролюватися, реєструватися та підлягати аудиту; у повідомленні про використання системи зазначено, що несанкціоноване використання системи заборонено і тягне за собою кримінальну та цивільну відповідальність; у повідомленні про використання системи зазначено, що використання системи означає згоду на моніторинг і запис дій; сповіщення або банер залишається на екрані до тих пір, поки користувачі не визнають умови використання та не приймуть явні дії для входу в систему або подальшого доступу до неї; для загальнодоступних систем інформація про використання системи умови відображається перед наданням подальшого доступу до загальнодоступної системи; для загальнодоступних систем відображаються будь-які посилання на моніторинг, запис або аудит, які узгоджуються з положеннями про конфіденційність таких систем, що зазвичай забороняють ці види діяльності; для загальнодоступних систем додається опис авторизованого використання системи = \textbf{default\_\allowbreak{}deny\_\allowbreak{}ruleabac\_\allowbreak{}rule\_\allowbreak{}1} \\ \hline
18 & БЛОКУВАННЯ ПРИСТРОЮ (AC-11) &  & AC-11 & Вибрано одне або декілька з наступних ЗНАЧЕНЬ ПАРАМЕТРІВ: \{ініціювання блокування пристрою після періоду неактивності; вимога до користувача ініціювати блокування пристрою перед тим, як залишити систему без нагляду\} = \textbf{30}\newline Часовий проміжок бездіяльності, після якого ініціюється блокування пристрою (якщо вибрано) = \textbf{30} \\ \hline
19 & id-spe-ac-11-1 &  & AC-11(1) &  \\ \hline
20 & Припинення сеансу (AC-12) &  & AC-12 & Визначено умови або події, що вимагають припинення сеансу = \textbf{loginlogoutfailed\_\allowbreak{}attempt} \\ \hline
21 & ВІДДАЛЕНИЙ ДОСТУП (AC-17) &  & AC-17 &  \\ \hline
22 & id-spe-ac-17-3 &  & AC-17(3) &  \\ \hline
23 & id-spe-ac-17-4 &  & AC-17(4) &  \\ \hline
24 & БЕЗДРОТОВИЙ ДОСТУП (AC-18) &  & AC-18 &  \\ \hline
25 & id-spe-ac-18-1 &  & AC-18(1) & Бездротовий доступ до системи захищений за допомогою шифрування = \textbf{AES-256-GCM} \\ \hline
26 & id-spe-ac-18-3 &  & AC-18(3) &  \\ \hline
27 & id-spe-ac-19 &  & AC-19 &  \\ \hline
28 & id-spe-ac-19-5 &  & AC-19(5) &  \\ \hline
29 & ВИКОРИСТАННЯ ЗОВНІШНІХ СИСТЕМ (AC-20) &  & AC-20 & Вибрано одне або більше з наступних ЗНАЧЕНЬ ПАРАМЕТРІВ: \{встановити умови та положення; визначити заходи захисту\} = \textbf{}\newline Визначено умови та положення, що відповідають довірчим відносинам, встановленим з іншими організаціями, які володіють, експлуатують та/або обслуговують зовнішні системи (якщо вибрано) = \textbf{} \\ \hline
30 & id-spe-ac-20-1 &  & AC-20(1) & Авторизовані особи мають право використовувати зовнішню систему для доступу до системи або для обробки, зберігання чи передачі інформації, що контролюється організацією, лише після перевірки виконання заходів безпеки та конфіденційності, зазначених у політиці безпеки та конфіденційності організації, а також планах безпеки та конфіденційності (якщо такі застосовуються) = \textbf{adminsecurity\_\allowbreak{}officer}\newline Авторизовані особи мають право використовувати зовнішню систему для доступу до системи або для обробки, зберігання чи передачі інформації, що контролюється організацією, лише після збереження погоджених угод про підключення або обробку системи з структурою організації, на якій розміщена зовнішня система = \textbf{adminsecurity\_\allowbreak{}officer} \\ \hline
31 & id-spe-ac-20-2 &  & AC-20(2) & Використання портативних пристроїв носіїв інформації уповноваженими особами обмежено у зовнішніх системах за допомогою обмеження  = \textbf{adminsecurity\_\allowbreak{}officer}\newline Визначено обмеження на використання авторизованими особами портативних носіїв інформації у зовнішніх системах = \textbf{adminsecurity\_\allowbreak{}officer} \\ \hline
32 & ПУБЛІЧНО ДОСТУПНИЙ КОНТЕНТ (AC-22) &  & AC-22 & Визначені особи уповноважені на розміщення інформації в загальнодоступній системі = \textbf{adminsecurity\_\allowbreak{}officer}\newline Уповноважені особи проходять навчання, щоб гарантувати, що загальнодоступна інформація не містить інформацію з обмеженим доступом = \textbf{adminsecurity\_\allowbreak{}officer}\newline Зміст у загальнодоступній системі перевіряється на наявність інформації з обмеженим доступом з частотою = \textbf{30}\newline Визначено частоту, з якою слід переглядати вміст загальнодоступної системи на предмет наявності там інформації з обмеженим доступом = \textbf{30} \\ \hline
\multicolumn{5}{|>{\raggedright\arraybackslash}p{14.5cm}|}{\textbf{2. ОБІЗНАНІСТЬ ТА НАВЧАННЯ (AT) (AT)}} \\ \hline
33 & id-spe-at-1 &  & AT-1 &  \\ \hline
34 & НАВЧАННЯ З ПІДВИЩЕННЯ ОБІЗНАНОСТІ (AT-2) &  & AT-2 & Оновлюється зміст навчання грамотності та підвищення обізнаності з частотою = \textbf{30}\newline Визначено періодичність проведення навчання грамотності з питань безпеки для користувачів системи (в тому числі менеджерів, вищого керівництва та підрядників) після початкового тренінгу = \textbf{щорічно}\newline Визначено періодичність проведення навчання грамотності з питань конфіденційності для користувачів системи (в тому числі менеджерів, вищого керівництва та підрядників) після початкового тренінгу = \textbf{щорічно}\newline Визначено події, які потребують навчання користувачів системи грамотності з питань безпеки = \textbf{loginlogoutfailed\_\allowbreak{}attempt}\newline Визначено частоту оновлення навчання грамотності та змісту обізнаності; AT-02\_\allowbreak{}ODP[07] визначено події після яких необхідне оновлення навчання грамотності та змісту обізнаності; AT-02(a)[01][01] навчання з грамотності з питань безпеки надається користувачам системи (включаючи менеджерів, керівників вищої ланки та підрядників) як частина початкового навчання для нових користувачів; AT-02(a)[01][02] навчання з грамотності з питань конфіденційності надається користувачам системи (включаючи менеджерів, керівників вищої ланки та підрядників) як частина початкового навчання для нових користувачів; AT-02(a)[01][03] для користувачів системи (включно з менеджерами, вищим керівництвом та підрядниками) проводиться навчання з безпеки з періодичністю після цього; AT-02(a)[01][04] для користувачів системи (включно з менеджерами, вищим керівництвом та підрядниками) проводиться навчання з конфіденційності з періодичністю після цього; AT-02(a)[02][01] тренінги з грамотності з питань безпеки проводяться для користувачів системи (включаючи менеджерів, керівників вищої ланки та підрядників), коли цього вимагають зміни в системі або після подій; AT-02(a)[02][02] тренінги з грамотності з питань конфіденціності проводяться для користувачів системи (включаючи менеджерів, керівників вищої ланки та підрядників), коли цього вимагають зміни в системі або після подій = \textbf{loginlogoutfailed\_\allowbreak{}attempt}\newline Визначено події, які потребують навчання користувачів системи грамотності з питань конфіденційності = \textbf{loginlogoutfailed\_\allowbreak{}attempt} \\ \hline
35 & id-spe-at-2-2 &  & AT-2(2) &  \\ \hline
36 & id-spe-at-2-3 &  & AT-2(3) &  \\ \hline
37 & РОЛЬОВЕ НАВЧАННЯ (AT-3) &  & AT-3 & Оновлюється вміст навчання на основі ролей з частотою = \textbf{30}\newline Визначено ролі та обов'язки для тренінгів з безпеки на основі ролей = \textbf{adminsecurity\_\allowbreak{}officer}\newline Визначено ролі та обов'язки для тренінгів з конфіденційності на основі ролей = \textbf{adminsecurity\_\allowbreak{}officer}\newline Визначено частоту проведення тренінгів на основі ролей з безпеки та конфіденційності для призначеного персоналу після початкової підготовки = \textbf{adminsecurity\_\allowbreak{}officer}\newline Визначено частоту оновлення змісту навчання на основі ролей = \textbf{30}\newline Визначено події, які потребують оновлення змісту навчання на основі ролей; AT-03(a)[01][01] навчання з безпеки на основі ролей проводиться для ролей та обов'язків перед авторизацією доступу до системи, інформації або виконанням призначених обов'язків; AT-03(a)[01][02] навчання з конфіденційності на основі ролей проводиться для ролей та обов'язків перед авторизацією доступу до системи, інформації або виконанням призначених обов'язків; AT-03(a)[01][03] навчання з безпеки на основі ролей проводиться для ролей та обов'язків з частотою після цього; AT-03(a)[01][04] навчання з конфіденційності на основі ролей проводиться для ролей та обов'язків з частотою після цього; AT-03(a)[02][01] навчання з питань безпеки на основі ролей проводиться для персоналу, який виконує певні ролі та обов'язки у сфері безпеки, коли цього вимагають зміни в системі; AT-03(a)[02][02] навчання з питань конфіденційності на основі ролей проводиться для персоналу, який виконує певні ролі та обов'язки у сфері безпеки, коли цього вимагають зміни в системі = \textbf{adminsecurity\_\allowbreak{}officer} \\ \hline
\multicolumn{5}{|>{\raggedright\arraybackslash}p{14.5cm}|}{\textbf{3. АУДИТ ТА ПІДЗВІТНІСТЬ (AU) (AU)}} \\ \hline
38 & id-spe-au-1 &  & AU-1 & Розроблено та задокументовано політику аудиту та підзвітності = \textbf{default\_\allowbreak{}deny\_\allowbreak{}ruleabac\_\allowbreak{}rule\_\allowbreak{}1}\newline Політика аудиту та підзвітності доведена до персонал або ролі = \textbf{adminsecurity\_\allowbreak{}officer}\newline Розроблені та задокументовані процедури аудиту та підзвітності, що сприяють впровадженню політики аудиту та підзвітності, а також відповідні заходи контролю аудиту та підзвітності = \textbf{default\_\allowbreak{}deny\_\allowbreak{}ruleabac\_\allowbreak{}rule\_\allowbreak{}1}\newline Посадова особа призначається для управління політикою та процедурами аудиту та підзвітності AU-01(c)[01][01] переглядається та оновлюється поточна політика аудиту та підзвітності з частота; AU-01(c)[01][02] переглядається та оновлюється поточна політика аудиту та підзвітності після подій; AU-01(c)[02][01] переглядається та оновлюється поточні процедури аудиту та підзвітності з частота; AU-01(c)[02][02] переглядається та оновлюється поточні процедури аудиту та підзвітності після подій = \textbf{adminsecurity\_\allowbreak{}officer}\newline Визначено персонал або ролі, до яких має бути доведена політика аудиту та підзвітності = \textbf{adminsecurity\_\allowbreak{}officer}\newline Визначено персонал або ролі, на які поширюються процедури аудиту та підзвітності = \textbf{adminsecurity\_\allowbreak{}officer}\newline Визначено посадову особу, яка управлятиме політикою та процедурами аудиту та підзвітності = \textbf{adminsecurity\_\allowbreak{}officer}\newline Визначено частоту, з якою переглядається та оновлюється поточна політика аудиту та підзвітності = \textbf{default\_\allowbreak{}deny\_\allowbreak{}ruleabac\_\allowbreak{}rule\_\allowbreak{}1}\newline Визначено події, які потребують перегляду та оновлення поточної політики аудиту та підзвітності = \textbf{default\_\allowbreak{}deny\_\allowbreak{}ruleabac\_\allowbreak{}rule\_\allowbreak{}1}\newline Визначено частоту, з якою переглядаються та оновлюються поточні процедури аудиту та підзвітності = \textbf{30}\newline Визначено події, які потребують перегляду та оновлення поточної процедури аудиту та підзвітності = \textbf{loginlogoutfailed\_\allowbreak{}attempt} \\ \hline
39 & ПОДІЇ АУДИТУ (AU-2) &  & AU-2 & Зазначені типи подій реєструються 02\_\allowbreak{}ODP[03] частота або ситуація>; <AU- = \textbf{щорічно}\newline Переглядаються та оновлюються типи подій, вибрані для реєстрації, частота. у системі = \textbf{щорічно}\newline Визначено частоту або ситуацію, що вимагає проведення аудиту для кожної ідентифікованої події = \textbf{loginlogoutfailed\_\allowbreak{}attempt}\newline Частота перегляду та оновлення типів подій, обраних для журналювання = \textbf{щорічно} \\ \hline
40 & ЗМІСТ ЗАПИСІВ АУДИТУ (AU-3) &  & AU-3 & Записи аудиту містять інформацію, яка встановлює який тип події стався = \textbf{loginlogoutfailed\_\allowbreak{}attempt}\newline Записи аудиту містять інформацію, яка встановлює джерело події = \textbf{loginlogoutfailed\_\allowbreak{}attempt}\newline Записи аудиту містять інформацію, яка встановлює наслідки події = \textbf{loginlogoutfailed\_\allowbreak{}attempt}\newline Записи аудиту містять інформацію, яка встановлює результат події та ідентифікатор будь-яких осіб або суб'єктів, пов'язаних з подією = \textbf{loginlogoutfailed\_\allowbreak{}attempt} \\ \hline
41 & id-spe-au-3-1 &  & AU-3(1) &  \\ \hline
42 & id-spe-au-5 &  & AU-5 & Персонал або ролі отримують сповіщення у разі збою процесу обробки даних аудиту періоду часу = \textbf{adminsecurity\_\allowbreak{}officer}\newline Додаткові дії виконуються у разі збою процесу обробки даних аудиту = \textbf{loginlogoutfailed\_\allowbreak{}attempt}\newline Визначено персонал або ролі, які отримують сповіщення про збої в процесі обробки даних аудиту  = \textbf{adminsecurity\_\allowbreak{}officer}\newline Визначено період часу, протягом якого персонал або ролі отримують сповіщення про збої в процесі обробки даних аудиту = \textbf{adminsecurity\_\allowbreak{}officer}\newline Визначено додаткові дії, яких слід вжити у випадку збою в процесі обробки даних аудиту  = \textbf{loginlogoutfailed\_\allowbreak{}attempt} \\ \hline
43 & ОГЛЯД, АНАЛІЗ І ЗВІТНІСТЬ АУДИТУ (AU-6) &  & AU-6 & Записи аудиту системи переглядаються та аналізуються частота для виявлення ознак неналежної або незвичної діяльності та потенційного впливу неналежної або незвичної діяльності = \textbf{щотижня}\newline Звіт аудиту відправляється персоналу або ролям = \textbf{adminsecurity\_\allowbreak{}officer}\newline Визначено частоту, з якою переглядаються та аналізуються записи аудиту системи = \textbf{30}\newline Визначено персонал або ролі які отримують результати оглядів та аналізів системних записів = \textbf{adminsecurity\_\allowbreak{}officer} \\ \hline
44 & id-spe-au-6-3 &  & AU-6(3) &  \\ \hline
45 & id-spe-au-7 &  & AU-7 & Забезпечено можливість скорочення записів перевірок аудитом та звітів, які не змінюють оригінальний зміст або час упорядкування записів аудиту = \textbf{30}\newline Реалізовано можливість скорочення записів перевірок аудитом та звітів, які не змінюють оригінальний зміст або час упорядкування записів аудиту = \textbf{30} \\ \hline
46 & ПОЗНАЧКА ЧАСУ (AU-8) &  & AU-8 & Внутрішній системний годинник використовується для створення позначок часу для записів аудиту = \textbf{30}\newline Позначки часу застосовуються для записів аудиту, які відповідають деталізація вимірювання часу і які використовують всесвітній координований час, мають фіксоване місцеве зміщення місцевого часу від всесвітнього координованого часу або включають зміщення місцевого часу як частину позначки часу = \textbf{30}\newline Визначено деталізацію вимірювання часу для часових позначок записів аудиту = \textbf{30} \\ \hline
47 & ЗАХИСТ ІНФОРМАЦІЇ АУДИТУ (AU-9) &  & AU-9 & Персонал або ролі отримують сповіщення при виявленні несанкціонованого доступу, зміни або видалення інформації аудиту = \textbf{adminsecurity\_\allowbreak{}officer}\newline Визначено персонал або ролі, які мають бути сповіщені при виявленні несанкціонованого доступу, зміни або видалення інформації аудиту = \textbf{adminsecurity\_\allowbreak{}officer} \\ \hline
48 & id-spe-au-9-4 &  & AU-9(4) &  \\ \hline
49 & ЗБЕРЕЖЕННЯ ЗАПИСІВ АУДИТУ (AU-11) &  & AU-11 & Записи аудиту зберігаються впродовж період часу, щоб забезпечити підтримку розслідування (постфактум) інцидентів безпеки та конфіденційності, а також відповідати нормативним та вимогам організації щодо збереження даних аудиту = \textbf{30}\newline Визначено період часу для зберігання записів аудиту, який узгоджується з політикою зберігання записів = \textbf{default\_\allowbreak{}deny\_\allowbreak{}ruleabac\_\allowbreak{}rule\_\allowbreak{}1} \\ \hline
50 & ГЕНЕРАЦІЯ ДАНИХ АУДИТУ (AU-12) &  & AU-12 & Персонал або ролі може/можуть вибирати типи подій, які будуть реєструватися певними компонентами системи; AU-12(c) згенеровано записи аудиту для типів подій, визначених у AU02\_\allowbreak{}ODP[02], які включають вміст записів аудиту, визначений у AU03 = \textbf{adminsecurity\_\allowbreak{}officer}\newline Визначено персонал або ролі, яким дозволено обирати типи подій, що мають реєструватися певними компонентами системи = \textbf{adminsecurity\_\allowbreak{}officer} \\ \hline
\multicolumn{5}{|>{\raggedright\arraybackslash}p{14.5cm}|}{\textbf{4. ОЦІНЮВАННЯ, АВТОРИЗАЦІЯ ТА МОНІТОРИНГ (CA) (CA)}} \\ \hline
51 & id-spe-ca-1 &  & CA-1 & Розроблено та задокументовано політику оцінювання, авторизації та моніторингу = \textbf{default\_\allowbreak{}deny\_\allowbreak{}ruleabac\_\allowbreak{}rule\_\allowbreak{}1}\newline Політика оцінювання, авторизації та моніторингу поширюється серед персоналу або ролей = \textbf{adminsecurity\_\allowbreak{}officer}\newline Розроблені та задокументовані процедури оцінювання, авторизації та моніторингу, що сприяють впровадженню політики оцінювання, авторизації та моніторингу, а також пов'язані з ними засоби контролю оцінювання, авторизації та моніторингу = \textbf{default\_\allowbreak{}deny\_\allowbreak{}ruleabac\_\allowbreak{}rule\_\allowbreak{}1}\newline Посадова особа призначається для управління розробкою, документуванням та розповсюдженням політики та процедур оцінювання, авторизації та моніторингу; CA-01(c)[01][01] переглядається та оновлюється поточна політика оцінки, авторизації та моніторингу з частотою; CA-01(c)[01][02] переглядається та оновлюється поточна політика оцінки, авторизації та моніторингу після подій; CA-01(c)[02][01] переглядаються та оновлюються поточні процедури оцінки, авторизації та моніторингу з частотою; CA-01(c)[02][02] переглядаються та оновлюються поточні процедури оцінки, авторизації та моніторингу після подій; ПОТЕНЦІЙНІ МЕТОДИ ТА ОБ'ЄКТИ ОЦІНКИ: = \textbf{adminsecurity\_\allowbreak{}officer}\newline Визначено персонал або ролі, серед яких має бути поширена політика оцінювання, авторизації та моніторингу = \textbf{adminsecurity\_\allowbreak{}officer}\newline Визначено персонал або ролі, серед яких мають бути поширені процедури оцінювання, авторизації та моніторингу = \textbf{adminsecurity\_\allowbreak{}officer}\newline Визначено посадову особу, яка управлятиме розробкою, документуванням і розповсюдженням політики та процедур оцінювання, авторизації та моніторингу = \textbf{adminsecurity\_\allowbreak{}officer}\newline Визначено частоту, з якою переглядається та оновлюється поточна політика оцінювання, авторизації та моніторингу = \textbf{default\_\allowbreak{}deny\_\allowbreak{}ruleabac\_\allowbreak{}rule\_\allowbreak{}1}\newline Визначено події, які потребують перегляду та оновлення поточної політики оцінки, авторизації та моніторингу = \textbf{default\_\allowbreak{}deny\_\allowbreak{}ruleabac\_\allowbreak{}rule\_\allowbreak{}1}\newline Визначено частоту, з якою переглядається та оновлюється поточні процедури оцінювання, авторизації та моніторингу = \textbf{30}\newline Визначено події, які потребують перегляду та оновлення поточні процедури оцінки, авторизації та моніторингу = \textbf{loginlogoutfailed\_\allowbreak{}attempt} \\ \hline
52 & ОЦІНЮВАННЯ (CA-2) &  & CA-2 & Розроблено план контрольної оцінки, який описує обсяг оцінки, включаючи заходи захисту та посилені заходи, що підлягають оцінюванню. CA-02(b)[02] розроблено план контрольної оцінки, який описує обсяг оцінки, включаючи процедури оцінювання, які використовуватимуться для визначення ефективності заходів. CA-02(b)[03][01] розроблено план контрольної оцінки, який описує обсяг оцінки, включаючи середовище оцінювання. CA-02(b)[03][02] розроблено план контрольної оцінки, який описує обсяг оцінки, включаючи групу оцінювання. CA-02(b)[03][03] розроблено план контрольної оцінки, який описує обсяг оцінки, включаючи ролі й обов'язки з оцінювання = \textbf{adminsecurity\_\allowbreak{}officer}\newline Розроблено план контрольної оцінки, який описує обсяг оцінки, включаючи процедури оцінювання, які використовуватимуться для визначення ефективності заходів. CA-02(b)[03][01] розроблено план контрольної оцінки, який описує обсяг оцінки, включаючи середовище оцінювання. CA-02(b)[03][02] розроблено план контрольної оцінки, який описує обсяг оцінки, включаючи групу оцінювання. CA-02(b)[03][03] розроблено план контрольної оцінки, який описує обсяг оцінки, включаючи ролі й обов'язки з оцінювання = \textbf{adminsecurity\_\allowbreak{}officer}\newline План оцінки заходів захисту розглядається та затверджується уповноваженою посадовою особою або призначеним представником перед проведенням оцінки = \textbf{adminsecurity\_\allowbreak{}officer}\newline Заходи захисту оцінюються в системі та середовищі її функціонування частота оцінки, щоб визначити, наскільки коректно реалізовані заходи захисту, чи працюють вони за призначенням і чи дають бажаний результат щодо дотримання встановлених вимог до безпеки = \textbf{щорічно}\newline Заходи захисту оцінюються в системі та середовищі її функціонування частота оцінки, щоб визначити, наскільки коректно реалізовані заходи захисту, чи працюють вони за призначенням і чи дають бажаний результат щодо дотримання встановлених вимог до конфіденційності; CA-02(e) готується звіт оцінювання = \textbf{щорічно}\newline Результати оцінювання з безпеки надаються особам або ролям. оцінювання , який документує результати = \textbf{adminsecurity\_\allowbreak{}officer}\newline Визначено частоту, з якою слід оцінювати засоби контролю в системі та середовищі її функціонування = \textbf{30}\newline Визначені особи або ролі, яким мають бути надані результати оцінювання з безпеки = \textbf{adminsecurity\_\allowbreak{}officer} \\ \hline
53 & ВЗАЄМОДІЯ СИСТЕМ (CA-3) &  & CA-3 & Характеристики інтерфейсу задокументовані як частина кожної угоди про обмін = \textbf{30}\newline Вимоги до безпеки задокументовані як частина кожної угоди про обмін = \textbf{30}\newline Вимоги щодо конфіденційності задокументовані як частина кожної угоди про обмін = \textbf{30}\newline Заходи захисту задокументовані як частина кожної угоди про обмін = \textbf{30}\newline Відповідальність за кожну систему задокументована як частина кожної угоди про обмін = \textbf{30}\newline Характер переданої інформації документується як частина кожної угоди про обмін = \textbf{30}\newline Угоди переглядаються та оновлюються частота = \textbf{щорічно}\newline Визначено частоту, з якою необхідно переглядати та оновлювати угоди = \textbf{30} \\ \hline
54 & id-spe-ca-5 &  & CA-5 & Розроблено план усунення недоліків та контрольні показники для системи, та задокументувано заплановані дії організації з коригування, спрямовані на усунення недоліків та зауважень, виявлених під час оцінки заходів захисту, а також на зменшення або усунення відомих вразливостей в системі = \textbf{loginlogoutfailed\_\allowbreak{}attempt}\newline Існуючий план оновлюються частота на основі результатів оцінювання заходів, незалежних аудитів та постійного моніторингу = \textbf{щорічно}\newline Визначено частоту оновлення чинного плану усунення недоліків та контрольних показників на основі результатів оцінювання заходів захисту, незалежних аудитів та постійного моніторингу = \textbf{30} \\ \hline
55 & БЕЗПЕРЕРВНИЙ МОНІТОРИНГ (CA-7) &  & CA-7 & Безперервний моніторинг на рівні системи включає встановлені частоти для моніторингу ефективності заходів захисту = \textbf{30}\newline Безперервний моніторинг на рівні системи включає встановлені частоти для оцінки ефективності заходів захисту = \textbf{30}\newline Безперервний моніторинг на рівні системи включає в себе дії з реагування на результати аналізу інформації, пов'язаної з безпекою та приватністю = \textbf{loginlogoutfailed\_\allowbreak{}attempt}\newline Безперервний моніторинг на рівні системи включає повідомлення про статус безпеки системи для персоналу або ролей частота = \textbf{adminsecurity\_\allowbreak{}officer}\newline Безперервний моніторинг на рівні системи включає повідомлення про стан конфіденційності системи персоналу або ролям частота = \textbf{adminsecurity\_\allowbreak{}officer}\newline Визначено частоту, з якою слід моніторити ефективність заходів захисту = \textbf{30}\newline Визначено частоту, з якою слід оцінювати ефективність заходів захисту = \textbf{30}\newline Визначено персонал або ролі, яким повідомляється про стан безпеки системи = \textbf{adminsecurity\_\allowbreak{}officer}\newline Визначено частоту, з якою повідомляється про стан безпеки системи = \textbf{30}\newline Визначено персонал або ролі, яким повідомляється про стан конфіденційності системи = \textbf{adminsecurity\_\allowbreak{}officer}\newline Визначено частоту, з якою повідомляється про стан конфіденційності системи = \textbf{30} \\ \hline
\multicolumn{5}{|>{\raggedright\arraybackslash}p{14.5cm}|}{\textbf{5. УПРАВЛІННЯ КОНФІГУРАЦІЄЮ (CM) (CM)}} \\ \hline
56 & id-spe-cm-1 &  & CM-1 & Розроблено та задокументовано політику управління конфігурацією = \textbf{default\_\allowbreak{}deny\_\allowbreak{}ruleabac\_\allowbreak{}rule\_\allowbreak{}1}\newline Політика управління конфігурацією поширюється на персонал або ролі = \textbf{adminsecurity\_\allowbreak{}officer}\newline Розроблені та задокументовані процедури управління , що сприяють реалізації політики управління конфігурацією та пов'язаних з нею заходів управління конфігурацією = \textbf{default\_\allowbreak{}deny\_\allowbreak{}ruleabac\_\allowbreak{}rule\_\allowbreak{}1}\newline Посадова особа призначається для управління розробкою, документуванням і розповсюдженням політики та процедур керування конфігурацією; CM-01(c)[01][01] переглядається та оновлюється поточна політика управління конфігурацією частота; CM-01(c)[01][02] переглядається та оновлюється поточна політика управління конфігурацією після події; CM-01(c)[02][01] переглядаються та оновлюються поточні процедури управління конфігурацією частота; CM-01(c)[02][02] переглядаються та оновлюються поточні процедури управління конфігурацією після події; ЗНАЧЕННЯ конфігурацією = \textbf{adminsecurity\_\allowbreak{}officer}\newline Визначено персонал або ролі, на яких поширюється політика управління конфігурацією = \textbf{adminsecurity\_\allowbreak{}officer}\newline Визначено персонал або ролі, на яких поширюється процедури управління конфігурацією = \textbf{adminsecurity\_\allowbreak{}officer}\newline Визначено посадову особу, яка управляє розробкою, документуванням і розповсюдженням політики та процедур керування конфігурацією = \textbf{adminsecurity\_\allowbreak{}officer}\newline Визначено частоту, з якою переглядається та оновлюється поточна політика управління конфігурацією = \textbf{default\_\allowbreak{}deny\_\allowbreak{}ruleabac\_\allowbreak{}rule\_\allowbreak{}1}\newline Визначено події, після яких переглядається та оновлюється поточна політика управління конфігурацією = \textbf{default\_\allowbreak{}deny\_\allowbreak{}ruleabac\_\allowbreak{}rule\_\allowbreak{}1}\newline Визначено частоту, з якою переглядаються та оновлюються поточні процедури управління конфігурацією = \textbf{30}\newline Визначено події, після яких переглядаються та оновлюються поточні процедури управління конфігурацією = \textbf{loginlogoutfailed\_\allowbreak{}attempt} \\ \hline
57 & БАЗОВА КОНФІГУРАЦІЯ (CM-2) &  & CM-2 & Переглядаються та оновлюються базові налаштування системи частота = \textbf{щорічно}\newline Визначено частоту налаштувань; перегляду та оновлення базових = \textbf{30} \\ \hline
58 & id-spe-cm-2-7 &  & CM-2(7) & Системи або компоненти системи з конфігураціями видаються особам, що перебувають у місцях, які організація вважає, становлять значний ризик = \textbf{adminsecurity\_\allowbreak{}officer}\newline Заходи безпеки застосовуються до систем або компонентів системи, коли особи повертаються з поїздки = \textbf{adminsecurity\_\allowbreak{}officer}\newline Визначено системи або компоненти систем, які мають видаватися особам, що перебувають у місцях зі значним ризиком = \textbf{adminsecurity\_\allowbreak{}officer} \\ \hline
59 & УПРАВЛІННЯ ЗМІНАМИ КОНФІГУРАЦІЇ (CM-3) &  & CM-3 & Записи про зміни конфігурації у системі зберігаються протягом <період часу CM-03\_\allowbreak{}ODP[01]> = \textbf{30}\newline Визначено період часу, протягом якого зберігатимуться записи про зміни конфігурації = \textbf{30}\newline Вибрано одне або декілька з наступних ЗНАЧЕНЬ ПАРАМЕТРІВ: \{частота; коли умови\} = \textbf{}\newline Визначено частоту, з якою викликаються елементи управління змінами конфігурації (якщо вибрано) = \textbf{30}\newline Визначено умови, за яких викликаються елементи управління змінами конфігурації (якщо вибрано); CM-03(a) визначено та задокументовано типи змін до системи, які контролюються конфігурацією = \textbf{} \\ \hline
60 & id-spe-cm-4 &  & CM-4 &  \\ \hline
61 & id-spe-cm-4-2 &  & CM-4(2) &  \\ \hline
62 & ОБМЕЖЕННЯ ДОСТУПУ ДО ЗМІНИ (CM-5) &  & CM-5 &  \\ \hline
63 & НАЛАШТУВАННЯ КОНФІГУРАЦІЇ (CM-6) &  & CM-6 & Зміни в налаштуваннях конфігурації відстежуються відповідно до політики та процедур організації = \textbf{default\_\allowbreak{}deny\_\allowbreak{}ruleabac\_\allowbreak{}rule\_\allowbreak{}1}\newline Зміни налаштувань конфігурації керуються відповідно до політики та процедур організації = \textbf{default\_\allowbreak{}deny\_\allowbreak{}ruleabac\_\allowbreak{}rule\_\allowbreak{}1} \\ \hline
64 & id-spe-cm-7 &  & CM-7 & Використання протоколи заборонено або обмежено = \textbf{TLS 1.3}\newline Визначено протоколи, які необхідно заборонити або обмежити; CM-07\_\allowbreak{}ODP[05] визначено програмне забезпечення, яке необхідно заборонити або обмежити = \textbf{TLS 1.3} \\ \hline
65 & id-spe-cm-7-1 &  & CM-7(1) & Система переглядається частота для виявлення непотрібних та/або незахищених функцій, портів, протоколів і послуг = \textbf{TLS 1.3}\newline Протоколи, які вважаються непотрібними та/або незахищеними, вимкнено = \textbf{TLS 1.3}\newline Визначено частоту, з якою слід переглядати систему для виявлення непотрібних та/або незахищених функцій, портів, протоколів і послуг = \textbf{TLS 1.3}\newline Визначено протоколи, які слід вимкнути, якщо вони вважаються непотрібними або незахищеними = \textbf{TLS 1.3} \\ \hline
66 & id-spe-cm-7-5 &  & CM-7(5) & Політика \textquotedbl{}заборонити все, дозволити за винятком\textquotedbl{} застосовуэться, щоб дозволити виконання авторизованих програм у системі = \textbf{default\_\allowbreak{}deny\_\allowbreak{}ruleabac\_\allowbreak{}rule\_\allowbreak{}1}\newline Переглядається та оновлюється список авторизованих програм частота = \textbf{}\newline Визначено частоту перегляду та оновлення списку авторизованих програм = \textbf{30} \\ \hline
67 & ІНВЕНТАРИЗАЦІЯ КОМПОНЕНТІВ СИСТЕМИ (CM-8) &  & CM-8 & Переглядається та оновлюється опис компонентів системи частота = \textbf{щорічно}\newline Визначено частоту перегляду та оновлення опису компонентів системи = \textbf{30} \\ \hline
68 & id-spe-cm-8-1 &  & CM-8(1) &  \\ \hline
69 & РОЗТАШУВАННЯ ІНФОРМАЦІЇ (CM-12) &  & CM-12 &  \\ \hline
\multicolumn{5}{|>{\raggedright\arraybackslash}p{14.5cm}|}{\textbf{6. ПЛАНУВАННЯ БЕЗПЕРЕРВНОЇ РОБОТИ (CP) (CP)}} \\ \hline
70 & РЕЗЕРВНЕ КОПІЮВАННЯ (CP-9) &  & CP-9 & Резервне копіювання інформації користувача, що міститься в компонентах системи, здійснюється частота = \textbf{щорічно}\newline Виконується резервне копіювання інформації системи, що міститься в системі частота; CP-09(с) створюються резервні копії документації системи, включаючи документацію, пов'язану з безпекою та конфіденційністю частота = \textbf{щорічно}\newline Визначено частоту, з якою слід проводити резервне копіювання інформації користувача відповідно до часу відновлення та цілей відновлення = \textbf{30}\newline Визначено частоту проведення резервного копіювання інформації системи, що відповідає завдань відновлення і встановлених цілей відновлення = \textbf{30}\newline Визначено частоту, з якою слід проводити резервне копіювання документації системи відповідно до часу відновлення та цілей точки відновлення = \textbf{30} \\ \hline
71 & id-spe-cp-9-8 &  & CP-9(8) & Реалізовано криптографічні механізми для запобігання несанкціонованому розкриттю та зміні резервної інформації = \textbf{AES-256-GCM} \\ \hline
\multicolumn{5}{|>{\raggedright\arraybackslash}p{14.5cm}|}{\textbf{7. ІДЕНТИФІКАЦІЯ ТА АВТЕНТИФІКАЦІЯ (IA) (IA)}} \\ \hline
72 & id-spe-ia-1 &  & IA-1 & Розроблено та задокументовано політику ідентифікації та автентифікації = \textbf{default\_\allowbreak{}deny\_\allowbreak{}ruleabac\_\allowbreak{}rule\_\allowbreak{}1}\newline Політика ідентифікації та автентифікації поширюється на персонал або ролі = \textbf{adminsecurity\_\allowbreak{}officer}\newline Розроблені та задокументовані процедури ідентифікації та автентифікації, що сприяють впровадженню політики ідентифікації та автентифікації, а також відповідні заходи ідентифікації та перевірки автентичності = \textbf{default\_\allowbreak{}deny\_\allowbreak{}ruleabac\_\allowbreak{}rule\_\allowbreak{}1}\newline Посадова особа призначається для управління політикою та процедурами ідентифікації та автентифікації; IA-01(c)[01][01] переглядається та оновлюється поточна політика ідентифікації та автентифікації частота; IA-01(c)[01][02] переглядається та оновлюється поточна політика ідентифікації та автентифікації після подій; IA-01(c)[02][01] переглядаються та оновлюються поточні процедури ідентифікації та автентифікації частота; IA-01(c)[02][02] переглядаються та оновлюються поточні процедури ідентифікації та автентифікації після подій = \textbf{adminsecurity\_\allowbreak{}officer}\newline Визначено персонал або ролі, на які поширюється політика ідентифікації та автентифікації = \textbf{adminsecurity\_\allowbreak{}officer}\newline Визначено персонал або ролі, на які поширюються процедури ідентифікації та автентифікації = \textbf{adminsecurity\_\allowbreak{}officer}\newline Визначено посадову особу, яка управляє політикою та процедурами ідентифікації та автентифікації = \textbf{adminsecurity\_\allowbreak{}officer}\newline Визначено частоту, з якою переглядається та оновлюється поточна політика ідентифікації та автентифікації = \textbf{default\_\allowbreak{}deny\_\allowbreak{}ruleabac\_\allowbreak{}rule\_\allowbreak{}1}\newline Визначено події, які потребують перегляду та оновлення поточної політики ідентифікації та автентифікації = \textbf{default\_\allowbreak{}deny\_\allowbreak{}ruleabac\_\allowbreak{}rule\_\allowbreak{}1}\newline Визначено частоту, з якою переглядаються та оновлюються поточні процедури ідентифікації та автентифікації = \textbf{30}\newline Визначено події, які потребують перегляду та оновлення процедур ідентифікації та автентифікації = \textbf{loginlogoutfailed\_\allowbreak{}attempt} \\ \hline
73 & id-spe-ia-2 &  & IA-2 &  \\ \hline
74 & id-spe-ia-2-1 &  & IA-2(1) &  \\ \hline
75 & id-spe-ia-2-2 &  & IA-2(2) &  \\ \hline
76 & id-spe-ia-2-8 &  & IA-2(8) &  \\ \hline
77 & id-spe-ia-3 &  & IA-3 &  \\ \hline
78 & УПРАВЛІННЯ ІДЕНТИФІКАЦІЄЮ (IA-4) &  & IA-4 & Управління ідентифікаторами здійснюється шляхом отримання дозволу від персоналу або ролей на призначення ідентифікатора особі, групі, ролі або пристрою = \textbf{adminsecurity\_\allowbreak{}officer}\newline Управління ідентифікаторами здійснюється шляхом вибору ідентифікатора, який ідентифікує окрему особу, групу, ролі або пристрій = \textbf{adminsecurity\_\allowbreak{}officer}\newline Управління ідентифікаторами здійснюється шляхом призначення ідентифікатора особі, групі, ролі або пристрою; IA-04(d) ідентифікатори управляються шляхом запобігання повторному використанню ідентифікаторів впродовж період часу = \textbf{adminsecurity\_\allowbreak{}officer}\newline Визначено персонал або ролі, від яких необхідно отримати дозвіл на призначення ідентифікатора = \textbf{adminsecurity\_\allowbreak{}officer}\newline Визначено період часу для запобігання повторному використанню ідентифікаторів = \textbf{30} \\ \hline
79 & id-spe-ia-4-4 &  & IA-4(4) &  \\ \hline
80 & АВТЕНТИФІКАТОР УПРАВЛІННЯ &  & IA-5 & Управління системними автентифікаторами здійснюється шляхом перевірки, як частини початкового розподілу автентифікатора, особи, групи, ролі або пристрою, який отримує автентифікатор = \textbf{adminsecurity\_\allowbreak{}officer}\newline Управління системними автентифікаторами здійснюється шляхом зміни/оновлення автентифікаторів у встановлений період часу або коли відбуваються події = \textbf{loginlogoutfailed\_\allowbreak{}attempt}\newline Визначено період часу для зміни або оновлення автентифікаторів за типом автентифікатора = \textbf{30} \\ \hline
81 & id-spe-ia-5-1 &  & IA-5(1) & Для автентифікації на основі паролів підтримується та оновлюється список часто використовуваних, очікуваних або скомпрометованих паролів частота, а також коли є підозра, що паролі організації були скомпрометовані прямо чи опосередковано = \textbf{adminsecurity\_\allowbreak{}officer}\newline Для автентифікації на основі паролів, коли паролі створюються або оновлюються користувачами, паролі перевіряються на відсутність у списку загальновживаних, очікуваних або скомпрометованих паролів в IA-05(01)(a) = \textbf{adminsecurity\_\allowbreak{}officer}\newline Для автентифікації на основі паролів, паролі передаються лише криптографічно захищеними каналами = \textbf{adminsecurity\_\allowbreak{}officer}\newline Для автентифікації на основі паролів паролі зберігаються за допомогою затвердженого алгоритму гешування, переважно використовуючи ключову геш-функцію = \textbf{adminsecurity\_\allowbreak{}officer}\newline Для автентифікації на основі пароля дозволяється вибір користувачем довгих паролів і фраз, що включають пробіли та всі друковані символи = \textbf{adminsecurity\_\allowbreak{}officer}\newline Для автентифікації на основі пароля використовуються автоматизовані інструменти, які допомагають користувачеві у виборі надійних автентифікаторів паролів = \textbf{adminsecurity\_\allowbreak{}officer}\newline Для автентифікації на основі пароля застосовуються склад та правила складності = \textbf{default\_\allowbreak{}deny\_\allowbreak{}ruleabac\_\allowbreak{}rule\_\allowbreak{}1}\newline Визначено частоту оновлення списку часто використовуваних, очікуваних або скомпрометованих паролів = \textbf{adminsecurity\_\allowbreak{}officer}\newline Визначено склад та правила складності автентифікатора = \textbf{default\_\allowbreak{}deny\_\allowbreak{}ruleabac\_\allowbreak{}rule\_\allowbreak{}1} \\ \hline
82 & id-spe-ia-6 &  & IA-6 & Забезпечино приховану зворотну передачу інформації автентифікації в про- цесі автентифікації для забезпечення захисту інформації від можливої експлуатації та використання неавторизованими особами = \textbf{adminsecurity\_\allowbreak{}officer} \\ \hline
83 & ПОВТОРНА АВТЕНТИФІКАЦІЯ (IA-11) &  & IA-11 &  \\ \hline
\multicolumn{5}{|>{\raggedright\arraybackslash}p{14.5cm}|}{\textbf{8. РЕАГУВАННЯ НА ІНЦИДЕНТИ (IR) (IR)}} \\ \hline
84 & id-spe-ir-1 &  & IR-1 & Розроблено та задокументовано політику реагування на інциденти = \textbf{default\_\allowbreak{}deny\_\allowbreak{}ruleabac\_\allowbreak{}rule\_\allowbreak{}1}\newline Політика реагування на інциденти поширюється серед персоналу або ролей = \textbf{adminsecurity\_\allowbreak{}officer}\newline Розроблені та задокументовані процедури реагування на інциденти, що сприяють впровадженню політики реагування на інциденти та пов'язаних з нею заходів захисту з реагування на інциденти = \textbf{default\_\allowbreak{}deny\_\allowbreak{}ruleabac\_\allowbreak{}rule\_\allowbreak{}1}\newline Посадова особа призначається для управління розробкою, документуванням та розповсюдженням політики та процедур реагування на інциденти; IR-01(c)[01][01] переглядається та оновлюється поточна політика реагування на інциденти частота; IR-01(c)[01][02] поточна політика реагування на інциденти переглядається та оновлюється після подій; IR-01(c)[02][01] переглядаються та оновлюються поточні процедури реагування на інциденти частота; IR-01(c)[02][02] поточні процедури реагування на інциденти переглядаються та оновлюються після подій = \textbf{adminsecurity\_\allowbreak{}officer}\newline Визначено персонал або ролі, до яких має бути доведена політика реагування на інциденти = \textbf{adminsecurity\_\allowbreak{}officer}\newline Визначено персонал або ролі, до яких мають бути доведені процедури реагування на інциденти = \textbf{adminsecurity\_\allowbreak{}officer}\newline Визначено посадову особу, яка керуватиме політикою та процедурами реагування на інциденти = \textbf{adminsecurity\_\allowbreak{}officer}\newline Визначено частоту, з якою переглядається та оновлюється поточна політика реагування на інциденти = \textbf{default\_\allowbreak{}deny\_\allowbreak{}ruleabac\_\allowbreak{}rule\_\allowbreak{}1}\newline Визначаються події, які потребують перегляду та оновлення поточної політики реагування на інциденти = \textbf{default\_\allowbreak{}deny\_\allowbreak{}ruleabac\_\allowbreak{}rule\_\allowbreak{}1}\newline Визначено частоту, з якою переглядаються та оновлюються поточні процедури реагування на інциденти = \textbf{30}\newline Визначено події, які потребують перегляду та оновлення процедур реагування на інциденти = \textbf{loginlogoutfailed\_\allowbreak{}attempt} \\ \hline
85 & НАВЧАННЯ З РЕАГУВАННЯ НА ІНЦИДЕНТИ (IR-2) &  & IR-2 & Навчання з реагування на інциденти надається користувачам системи відповідно до призначених ролей та обов'язків протягом періоду часу з моменту прийняття на себе ролі або обов'язків з реагування на інциденти або отримання доступу до системи = \textbf{adminsecurity\_\allowbreak{}officer}\newline Користувачам системи надається навчання з реагування на інциденти відповідно до призначених ролей та обов'язків частота = \textbf{щорічно}\newline Зміст навчання з реагування на інциденти переглядається та оновлюється частота = \textbf{щорічно}\newline Визначено частоту, з якою користувачі повинні проходити навчання з реагування на інциденти = \textbf{30}\newline Визначено частоту перегляду та оновлення змісту навчання з реагування на інциденти = \textbf{30}\newline Визначено події, які ініціюють перегляд змісту навчання з реагування на інциденти = \textbf{loginlogoutfailed\_\allowbreak{}attempt} \\ \hline
86 & ПЕРЕВІРКА РЕАГУВАНЬ НА ІНЦИДЕНТИ (IR-3) &  & IR-3 & Ефективність реагування системи на інциденти перевіряється частота за допомогою тестів = \textbf{щорічно}\newline Визначено частоту, з якою необхідно перевіряти ефективність реагування системи на інциденти = \textbf{30} \\ \hline
87 & Обробка інциденту (IR-4) &  & IR-4 &  \\ \hline
88 & Моніторинг інциденту (IR-5) &  & IR-5 &  \\ \hline
89 & ЗВІТНІСТЬ ПРО ІНЦИДЕНТИ (IR-6) &  & IR-6 & Персонал зобов'язаний повідомляти про підозрілі інциденти протягом періоду часу = \textbf{adminsecurity\_\allowbreak{}officer}\newline Визначено період часу, протягом якого персонал повинен повідомляти про підозрілі інциденти до уповноваженого органу = \textbf{adminsecurity\_\allowbreak{}officer} \\ \hline
90 & ПІДТРИМКА РЕАГУВАННЯ НА ІНЦИДЕНТИ (IR-7) &  & IR-7 &  \\ \hline
91 & План реагування на інциденти (IR-8) &  & IR-8 & Розроблено план реагування на інцидент, який розглядається та затверджується персоналом або ролями частота = \textbf{adminsecurity\_\allowbreak{}officer}\newline Розроблено план реагування на інциденти, в якому чітко визначено відповідальність за реагування на інциденти для організацій, персоналу або ролей = \textbf{adminsecurity\_\allowbreak{}officer}\newline Копії плану реагування на інцидент розповсюджуються серед <IR- 08\_\allowbreak{}ODP[04] персоналу з реагування на інциденти> = \textbf{adminsecurity\_\allowbreak{}officer}\newline План реагування на інциденти оновлюється з урахуванням змін у системі та організації або проблем, що виникають під час впровадження, виконання або тестування плану = \textbf{30}\newline Зміни в плані реагування на інцидент повідомляються персоналу з реагування на інциденти = \textbf{adminsecurity\_\allowbreak{}officer}\newline Визначено персонал або ролі, які переглядають та затверджують план реагування на інциденти = \textbf{adminsecurity\_\allowbreak{}officer}\newline Визначено періодичність перегляду та затвердження плану реагування на інциденти = \textbf{щорічно}\newline Визначені організації, персонал або ролі, які несуть відповідальність за реагування на інциденти; IR-08\_\allowbreak{}ODP[04] визначено персонал з реагування на інцидент (ідентифікований за іменами та/або за ролями), якому мають бути роздані копії плану реагування на інцидент = \textbf{adminsecurity\_\allowbreak{}officer}\newline Визначено персонал з реагування на інцидент (ідентифікований за іменами та/або ролями), якому повідомляються зміни до плану реагування на інцидент = \textbf{adminsecurity\_\allowbreak{}officer} \\ \hline
\multicolumn{5}{|>{\raggedright\arraybackslash}p{14.5cm}|}{\textbf{9. ТЕХНІЧНЕ ОБСЛУГОВУВАННЯ (MA) (MA)}} \\ \hline
92 & id-spe-ma-1 &  & MA-1 & Розроблено та задокументовано політику технічного обслуговування = \textbf{default\_\allowbreak{}deny\_\allowbreak{}ruleabac\_\allowbreak{}rule\_\allowbreak{}1}\newline Політика технічного обслуговування поширюється на персонал або ролі = \textbf{adminsecurity\_\allowbreak{}officer}\newline Розроблені та задокументовані процедури технічного обслуговування, що сприяють впровадженню політики технічного обслуговування та пов'язаних з нею заходів технічного обслуговування = \textbf{default\_\allowbreak{}deny\_\allowbreak{}ruleabac\_\allowbreak{}rule\_\allowbreak{}1}\newline Посадова особа призначається для управління розробкою, документуванням та розповсюдженням політики та процедур технічного обслуговування; MA-01(c)[01][01] переглядається та оновлюється поточна політика обслуговування частота; MA-01(c)[01][02] переглядається та оновлюється поточна політика обслуговування після подій; MA-01(c)[02][01] переглядаються та оновлюються поточні процедури технічного обслуговування частота; MA-01(c)[02][02] переглядаються та оновлюються поточні процедури технічного обслуговування після подій = \textbf{adminsecurity\_\allowbreak{}officer}\newline Визначено персонал або ролі, на які поширюється політика технічного обслуговування = \textbf{adminsecurity\_\allowbreak{}officer}\newline Визначено персонал або ролі, на які поширюються процедури технічного обслуговування = \textbf{adminsecurity\_\allowbreak{}officer}\newline Визначено посадову особу, яка керуватиме політикою та процедурами технічного обслуговування = \textbf{adminsecurity\_\allowbreak{}officer}\newline Визначено частоту, з якою переглядається та оновлюється поточна політика технічного обслуговування = \textbf{default\_\allowbreak{}deny\_\allowbreak{}ruleabac\_\allowbreak{}rule\_\allowbreak{}1}\newline Визначено події, які потребують перегляду та оновлення поточної політики технічного обслуговування = \textbf{default\_\allowbreak{}deny\_\allowbreak{}ruleabac\_\allowbreak{}rule\_\allowbreak{}1}\newline Визначено частоту, з якою переглядаються та оновлюються поточні процедури технічного обслуговування = \textbf{30}\newline Визначено події, які потребують перегляду та оновлення процедур технічного обслуговування = \textbf{loginlogoutfailed\_\allowbreak{}attempt} \\ \hline
93 & ІНСТРУМЕНТИ ДЛЯ ОБСЛУГОВУВАННЯ (MA-3) &  & MA-3 & Переглядаються раніше затверджені інструменти технічного обслуговування частота = \textbf{щорічно}\newline Визначено частоту, з якою слід переглядати раніше затверджені інструменти технічного обслуговування = \textbf{30} \\ \hline
94 & id-spe-ma-3-1 &  & MA-3(1) & Оглядаються інструменти для технічного обслуговування, які доставлені на об'єкт обслуговуючим персоналом, на предмет неправильних або несанкціонованих модифікацій = \textbf{adminsecurity\_\allowbreak{}officer} \\ \hline
95 & id-spe-ma-3-2 &  & MA-3(2) &  \\ \hline
96 & id-spe-ma-3-3 &  & MA-3(3) & Переміщення обладнання для технічного обслуговування, що містить інформацію організації, запобігається шляхом отримання дозволу від персонал або ролі = \textbf{adminsecurity\_\allowbreak{}officer}\newline Визначено персонал або ролі, які можуть надавати дозвіл на переміщення обладнання з об'єкту = \textbf{adminsecurity\_\allowbreak{}officer} \\ \hline
97 & ВІДДАЛЕНЕ ОБСЛУГОВУВАННЯ (MA-4) &  & MA-4 & Впроваджено віддалені дії з обслуговування та діагностики = \textbf{loginlogoutfailed\_\allowbreak{}attempt}\newline Відстежуються віддалені дії з обслуговування та діагностики = \textbf{loginlogoutfailed\_\allowbreak{}attempt}\newline Використання віддалених засобів технічного обслуговування та діагностики дозволено лише відповідно до політики організації = \textbf{default\_\allowbreak{}deny\_\allowbreak{}ruleabac\_\allowbreak{}rule\_\allowbreak{}1} \\ \hline
98 & ТЕХНІЧНИЙ ПЕРСОНАЛ (MA-5) &  & MA-5 & Запроваджено процес авторизації технічного персоналу = \textbf{adminsecurity\_\allowbreak{}officer}\newline Ведеться перелік авторизованих організацій або персоналу з технічного обслуговування = \textbf{adminsecurity\_\allowbreak{}officer}\newline Персонал без супроводу, який виконує технічне обслуговування системи, має необхідні дозволи на доступ = \textbf{adminsecurity\_\allowbreak{}officer}\newline Персонал організації з необхідними повноваженнями доступу та технічною компетентністю призначений/призначені для нагляду за діяльністю з технічного обслуговування персоналу, який не має необхідних дозволів на доступ = \textbf{adminsecurity\_\allowbreak{}officer} \\ \hline
\multicolumn{5}{|>{\raggedright\arraybackslash}p{14.5cm}|}{\textbf{10. ЗАХИСТ НОСІЇВ ІНФОРМАЦІЇ (MP) (MP)}} \\ \hline
99 & id-spe-mp-1 &  & MP-1 & Розроблено та задокументовано політику захисту носіїв інформації = \textbf{default\_\allowbreak{}deny\_\allowbreak{}ruleabac\_\allowbreak{}rule\_\allowbreak{}1}\newline Політика захисту носіїв інформації поширюється на персонал або ролі = \textbf{adminsecurity\_\allowbreak{}officer}\newline Розроблено та задокументовано процедури захисту носіїв інформації, що сприятимуть реалізації політики захисту носіїв інформації та заходів захисту носіїв інформації = \textbf{default\_\allowbreak{}deny\_\allowbreak{}ruleabac\_\allowbreak{}rule\_\allowbreak{}1}\newline Посадова особа призначається для управління розробкою, документуванням та розповсюдженням політики та процедур захисту носіїв інформації. MP-01(c)[01][01] переглядається та оновлюється поточна політика захисту носіїв інформації частота; MP-01(c)[01][02] переглядається та оновлюється поточна політика захисту носіїв інформації після подій; MP-01(c)[02][01] переглядаються та оновлюються поточні процедури захисту носіїв інформації частота; MP-01(c)[02][02] переглядаються та оновлюються поточні процедури захисту носіїв інформації після подій = \textbf{adminsecurity\_\allowbreak{}officer}\newline Визначено персонал або ролі, серед яких має бути поширена політика захисту носіїв інформації = \textbf{adminsecurity\_\allowbreak{}officer}\newline Визначено персонал або ролі, серед яких мають бути поширені процедури захисту носіїв інформації = \textbf{adminsecurity\_\allowbreak{}officer}\newline Визначено посадову особу, яка керуватиме політикою та процедурами захисту носіїв інформації = \textbf{adminsecurity\_\allowbreak{}officer}\newline Визначено частоту, з якою переглядається та оновлюється поточна політика захисту носіїв інформації = \textbf{default\_\allowbreak{}deny\_\allowbreak{}ruleabac\_\allowbreak{}rule\_\allowbreak{}1}\newline Визначено події, які потребують перегляду та оновлення чинної політики захисту носіїв інформації = \textbf{default\_\allowbreak{}deny\_\allowbreak{}ruleabac\_\allowbreak{}rule\_\allowbreak{}1}\newline Визначено частоту, з якою переглядаються та оновлюються чинні процедури захисту носіїв інформації = \textbf{30}\newline Визначено події, які потребують перегляду та оновлення процедур захисту носіїв інформації = \textbf{loginlogoutfailed\_\allowbreak{}attempt} \\ \hline
100 & ДОСТУП ДО НОСІЇВ ІНФОРМАЦІЇ (MP-2) &  & MP-2 & Доступ до типів цифрових носіїв інформації обмежено для персоналу або ролей = \textbf{adminsecurity\_\allowbreak{}officer}\newline Доступ до типів нецифрових носіїв інформації обмежено для персоналу або ролей = \textbf{adminsecurity\_\allowbreak{}officer}\newline Визначено персонал або ролі, уповноважені на доступ до цифрових носіїв інформації = \textbf{adminsecurity\_\allowbreak{}officer}\newline Визначено персонал або ролі, уповноважені на доступ до нецифрових носіїв інформації = \textbf{adminsecurity\_\allowbreak{}officer} \\ \hline
101 & МАРКУВАННЯ НОСІЇВ ІНФОРМАЦІЇ (MP-3) &  & MP-3 & Визначено типи носіїв інформації, які звільняються від маркування під час перебування на контрольованих зонах = \textbf{30} \\ \hline
102 & ЗБЕРІГАННЯ НОСІЇВ ІНФОРМАЦІЇ (MP-4) &  & MP-4 &  \\ \hline
103 & Транспортування носіїв інформації (MP-5) &  & MP-5 & Типи носіїв інформації системи захищаються під час транспортування за межі контрольованих зон за допомогою заходів безпеки = \textbf{30}\newline Типи носіїв інформації системи контролюються під час транспортування за межі контрольованих зон за допомогою заходів безпеки = \textbf{30}\newline Під час транспортування за межі контрольованих зон ведеться облік носіїв інформації системи = \textbf{30}\newline Визначено персонал, уповноважений здійснювати діяльність з транспортування носіїв інформації = \textbf{adminsecurity\_\allowbreak{}officer}\newline Діяльність, пов'язана з транспортуванням носіїв інформації системи, обмежується визначеним уповноваженим персоналом = \textbf{adminsecurity\_\allowbreak{}officer}\newline Визначено типи носіїв інформації системи для захисту та контролю під час транспортування за межі контрольованих зон = \textbf{30} \\ \hline
104 & id-spe-mp-6 &  & MP-6 & Застосовуються механізми очищення, надійність і цілісність яких відповідає категорії безпеки або рівню секретності інформації = \textbf{автоматизований засіб моніторингу} \\ \hline
105 & ВИКОРИСТАННЯ НОСІЇВ ІНФОРМАЦІЇ (MP-7) &  & MP-7 &  \\ \hline
\multicolumn{5}{|>{\raggedright\arraybackslash}p{14.5cm}|}{\textbf{11. id-spe-pe (PE)}} \\ \hline
106 & id-spe-pe-1 &  & PE-1 & Розроблено та задокументовано політику фізичного захисту та захисту робочого середовища = \textbf{default\_\allowbreak{}deny\_\allowbreak{}ruleabac\_\allowbreak{}rule\_\allowbreak{}1}\newline Політика фізичного захисту та захисту робочого середовища розповсюджується серед персоналу або ролей = \textbf{adminsecurity\_\allowbreak{}officer}\newline Посадова особа призначається для управління розробкою, документуванням та розповсюдженням політики та процедур фізичного захисту та захисту робочого середовища; PE-01(c)[01][01] переглядається та оновлюється поточна політика фізичного захисту та захисту робочого середовища частота; PE-01(c)[01][02] поточна політика фізичного захисту та захисту робочого середовища переглядається та оновлюється після подій; PE-01(c)[02][01] переглядаються та оновлюються поточні процедури фізичного захисту та захисту робочого середовища частота; поточні процедури фізичного та екологічного захисту переглядаються та оновлюються після подій. PE-01(c)[02][02] = \textbf{adminsecurity\_\allowbreak{}officer}\newline Визначено персонал або ролі, до яких має бути доведена політика фізичного захисту та захисту робочого середовища = \textbf{adminsecurity\_\allowbreak{}officer}\newline Визначено персонал або ролі, на які поширюються процедури фізичного захисту та захисту робочого середовища = \textbf{adminsecurity\_\allowbreak{}officer}\newline Визначено посадову особу, яка керуватиме політикою та процедурами фізичного захисту та захисту робочого середовища = \textbf{adminsecurity\_\allowbreak{}officer}\newline Визначено частоту, з якою переглядається та оновлюється поточна політика фізичного захисту та захисту робочого середовища = \textbf{default\_\allowbreak{}deny\_\allowbreak{}ruleabac\_\allowbreak{}rule\_\allowbreak{}1}\newline Визначено події, які потребують перегляду та оновлення поточної політики фізичного захисту та захисту робочого середовища = \textbf{default\_\allowbreak{}deny\_\allowbreak{}ruleabac\_\allowbreak{}rule\_\allowbreak{}1}\newline Визначено частоту, з якою переглядаються та оновлюються поточні процедури фізичного захисту та захисту робочого середовища = \textbf{30}\newline Визначено події, які потребують перегляду та оновлення процедур фізичного захисту та захисту робочого середовища = \textbf{loginlogoutfailed\_\allowbreak{}attempt} \\ \hline
107 & РЕ-2 доступу (PE-2) &  & PE-2 & Розроблено перелік осіб, які мають право авторизованого доступу до об'єкта, де перебуває система = \textbf{}\newline Затверджено перелік осіб, які мають право авторизованого доступу до об'єкта, де перебуває система = \textbf{}\newline Ведеться перелік осіб, які мають право авторизованого доступу до об'єкта, де перебуває система = \textbf{}\newline Переглядається список доступу, , у якому закріплений перелік персоналу або ролей, яким дозволений санкціонований доступ до об'єкта частота = \textbf{adminsecurity\_\allowbreak{}officer}\newline Особи видаляються зі списку доступу до об'єкта, коли доступ більше не потрібен = \textbf{adminsecurity\_\allowbreak{}officer}\newline Визначено періодичність перегляду списку доступу, у якому закріплений перелік персоналу або ролей, яким дозволений санкціонований доступ до об'єкта = \textbf{adminsecurity\_\allowbreak{}officer} \\ \hline
108 & ФІЗИЧНИЙ ДОСТУП ДО СИСТЕМИ &  & PE-3 & Активність відвідувачів контролюється умови = \textbf{}\newline Пристрої фізичного доступу інвентаризуються частота = \textbf{щорічно}\newline Коди доступу змінюється частота, коли код скомпрометовано, або коли особи, які володіють кодом, переводяться або звільняються; <PE- PE-03(g)[02] ключі змінюються частота, коли ключі втрачено, або коли особи, що володіють ключами, переводяться або звільняються = \textbf{adminsecurity\_\allowbreak{}officer}\newline Визначено умови, що вимагають супроводу відвідувачів та моніторингу активності відвідувачів = \textbf{}\newline Визначено частоту проведення інвентаризації пристроїв фізичного доступу = \textbf{30}\newline Визначено частоту, з якою потрібно змінювати коди доступу = \textbf{30}\newline Визначено частоту, з якою потрібно змінювати ключі = \textbf{30} \\ \hline
109 & РЕ-4 ліній електроживлення (PE-4) &  & PE-4 &  \\ \hline
110 & id-spe-pe-5 &  & PE-5 &  \\ \hline
111 & МОНІТОРИНГ ФІЗИЧНОГО ДОСТУПУ &  & PE-6 & Переглядаються журнали фізичного доступу частота = \textbf{щорічно}\newline Визначено частоту перегляду журналів фізичного доступу = \textbf{30}\newline Визначено події або потенційні ознаки подій, що вимагають перегляду журналів фізичного доступу = \textbf{loginlogoutfailed\_\allowbreak{}attempt} \\ \hline
112 & АЛЬТЕРНАТИВНЕ РОБОЧЕ МІСЦЕ (PE-17) &  & PE-17 & Визначаються заходи захисту, які будуть застосовуватися на альтернативних робочих місцях; РЕ-17(a) альтернативні робочі місця визначені та задокументовані; РЕ-17(b) заходи захисту впроваджені на альтернативних робочих місцях; РЕ-17(c) оцінюється ефективність заходів захисту на альтернативних робочих місцях; РЕ-17(d) працівникам надаються засоби комунікації з персоналом служби інформаційної безпеки на випадок інцидентів = \textbf{adminsecurity\_\allowbreak{}officer} \\ \hline
\multicolumn{5}{|>{\raggedright\arraybackslash}p{14.5cm}|}{\textbf{12. ПЛАНУВАННЯ (PL) (PL)}} \\ \hline
113 & id-spe-pl-1 &  & PL-1 & Визначено персонал або ролі, до яких має бути доведена політика планування безпеки = \textbf{adminsecurity\_\allowbreak{}officer}\newline Визначено персонал або ролі, на які поширюватимуться процедури планування безпеки = \textbf{adminsecurity\_\allowbreak{}officer}\newline Визначено посадову особу, яка керуватиме політикою та процедурами планування безпеки = \textbf{adminsecurity\_\allowbreak{}officer}\newline Визначено періодичність перегляду та оновлення поточної політики планування безпеки = \textbf{default\_\allowbreak{}deny\_\allowbreak{}ruleabac\_\allowbreak{}rule\_\allowbreak{}1}\newline Є події, які потребують перегляду та оновлення поточної політики планування безпеки = \textbf{default\_\allowbreak{}deny\_\allowbreak{}ruleabac\_\allowbreak{}rule\_\allowbreak{}1}\newline Визначена частота, з якою переглядаються та оновлюються поточні процедури планування безпеки = \textbf{щорічно} \\ \hline
114 & id-spe-pl-2 &  & PL-2 & Призначені особи або групи, з якими пов'язана діяльність з безпекою та конфіденційністю, що впливає на систему, яка потребує планування та координації = \textbf{adminsecurity\_\allowbreak{}officer}\newline Призначено персонал або ролі для отримання розповсюджуваних копій планів захисту інформації та конфіденційності системи = \textbf{adminsecurity\_\allowbreak{}officer}\newline Визначено періодичність перегляду інформації та конфіденційності системи; планів захисту PL-02a.01[01] розроблено план захисту інформації, архітектурі підприємства організації; який відповідає PL-02a.01[02] розроблено план конфіденційності для системи, відповідає архітектурі підприємства організації; який PL-02a.02[01] розроблено план захисту інформації, який чітко визначає складові компоненти системи; PL-02a.02[02] розроблено для системи план забезпечення конфіденційності, який чітко визначає складові компоненти системи; PL-02a.03[01] розроблено план захисту інформації системи, який описує операційний контекст системи з точки зору місії та бізнеспроцесів; PL-02a.03[02] розроблено для системи план забезпечення конфіденційності, який описує операційний контекст системи з точки зору місії та бізнес-процесів; PL-02a.04[01] розроблено план захисту інформації, який визначає осіб, що виконують системні ролі та обов'язки; PL-02a.04[02] розроблено для системи план забезпечення конфіденційності, який визначає осіб, що виконують ролі та обов'язки в системі; PL-02a.05[01] розроблено план захисту інформації, який визначає типи інформації, що обробляється, зберігається та передається системою; PL-02a.05[02] розроблено план конфіденційності для системи, який визначає типи інформації, що обробляється, зберігається та передається системою; PL-02a.06[01] розроблено план захисту інформації, який передбачає категоризацію безпеки системи, включаючи відповідне обґрунтування; PL-02a.06[02] розроблено для системи план забезпечення конфіденційності, який передбачає категоризацію системи за рівнем безпеки, включаючи обґрунтування; PL-02a.07[01] розроблено план захисту інформації, який описує будь-які конкретні загрози для системи, що викликають потенційні ризики для рганізації; PL-02a.07[02] розроблено план конфіденційності для системи, який описує будь-які конкретні загрози для системи, що викликають потенційні ризики для організації; PL-02a.08[01] розроблено план захисту інформації, який містить результати оцінки ризиків конфіденційності для систем, що обробляють інформацію, яка ідентифікує особу; PL-02a.08[02] розроблено для системи план забезпечення конфіденційності, який містить результати оцінки ризиків конфіденційності для систем, що обробляють інформацію, яка ідентифікує особу; PL-02a.09[01] розроблено план захисту інформації, який описує операційне середовище системи та будь-які залежності або зв'язки з іншими системами чи компонентами системи; PL-02a.09[02] розроблено для системи план забезпечення конфіденційності, який описує робоче середовище системи та будь-які залежності або зв'язки з іншими системами чи компонентами системи; PL-02a.10[01] розроблено план захисту інформації, який містить огляд вимог до безпеки системи; PL-02a.10[02] розроблено для системи план забезпечення конфіденційності, який містить огляд вимог до конфіденційності системи; PL-02a.11[01] розроблено план захисту інформації, який визначає будь-які відповідні базові рівні контролю або обмеження, якщо такі є; PL-02a.11[02] розроблено для системи план забезпечення конфіденційності, який визначає будь-які відповідні базові рівні контролю або обмеження, якщо такі є; PL-02a.12[01] розроблено план захисту інформації, який описує наявні або заплановані засоби контролю для виконання вимог безпеки, включаючи обґрунтування будь-яких рішень щодо адаптації; PL-02a.12[02] розроблено для системи план забезпечення конфіденційності, який описує наявні або заплановані засоби контролю для виконання вимог щодо конфіденційності, включаючи обґрунтування будь-яких рішень, пов'язаних з адаптацією; PL-02a.13[01] розроблено план захисту інформації, який включає визначення ризиків для архітектури безпеки та проектних рішень; PL-02a.13[02] розроблено план забезпечення конфіденційності для системи, який включає визначення ризиків для архітектури конфіденційності та проектних рішень; PL-02a.14[01] розроблено захисту інформації, який включає діяльність, пов'язану з безпекою, що впливає на систему і потребує планування та координації з окремими особами або групами; PL-02a.14[02] розроблено для системи план забезпечення конфіденційності, який включає діяльність, пов'язану з конфіденційністю, що впливає на систему і потребує планування та координації з окремими особами або групами; PL-02a.15[01] розроблено план захисту інформації, який розглядається та затверджується уповноваженою посадовою особою або призначеним представником до початку реалізації плану; PL-02a.15[02] розроблено план забезпечення конфіденційності для системи, який перевіряється та затверджується уповноваженою посадовою особою або призначеним представником перед впровадженням плану. PL-02b.[01] розповсюджуються копії персонал або ролі>; PL-02b.[02] повідомляються наступні зміни до планів персонал або ролі; PL-02c. переглядаються плани частота; планів серед <PL-02\_\allowbreak{}ODP[02] PL-02d.[01] оновлюються плани відповідно до змін у системі та середовищі діяльності; PL-02d.[02] оновлюються плани для вирішення проблем, виявлених під час реалізації плану; PL-02d.[03] оновлюються плани для вирішення проблем, виявлених під час контрольних оцінок; PL-02e.[01] захищені плани від несанкціонованого розголошення; PL-02e.[02] захищені плани від несанкціонованої модифікації = \textbf{adminsecurity\_\allowbreak{}officer} \\ \hline
115 & ПРАВИЛА ПОВЕДІНКИ (PL-4) &  & PL-4 & Визначено періодичність перегляду та оновлення правил поведінки = \textbf{щорічно}\newline Вибрано одне або декілька з наступних ЗНАЧЕНЬ ПАРАМЕТРА:\{частота; коли правила переглядаються або оновлюються\} = \textbf{default\_\allowbreak{}deny\_\allowbreak{}ruleabac\_\allowbreak{}rule\_\allowbreak{}1}\newline Визначена періодичність перегляду та повторного підтвердження правил поведінки (якщо вибрано); PL-04a.[01] встановлені правила, які описують обов'язки та очікувану поведінку щодо використання інформації та системи, безпеки та конфіденційності для осіб, яким потрібен доступ до системи; PL-04a.[02] надаються правила, які описують обов'язки та очікувану поведінку щодо використання інформації та системи, безпеки та конфіденційності, особам, які отримують доступ до системи; PL-04b. отримано перед наданням доступу до інформації та системи задокументоване підтвердження від таких осіб про те, що вони прочитали, зрозуміли та згодні дотримуватися правил поведінки; PL-04c. переглядаються та оновлюються правила поведінки частота; PL-04d. потрібно особам, які визнали попередню версію правил поведінки, прочитати та повторно визнати ЗНАЧЕННЯ ВИБРАНОГО ПАРАМЕТРА(ів) = \textbf{adminsecurity\_\allowbreak{}officer} \\ \hline
\multicolumn{5}{|>{\raggedright\arraybackslash}p{14.5cm}|}{\textbf{13. БЕЗПЕКА ПЕРСОНАЛУ (PS) (PS)}} \\ \hline
116 & id-spe-ps-1 &  & PS-1 & Визначено персонал або ролі, на які поширюється політика кадрової безпеки = \textbf{adminsecurity\_\allowbreak{}officer}\newline Визначено персонал або ролі, на які поширюються процедури кадрової безпеки = \textbf{adminsecurity\_\allowbreak{}officer}\newline Визначено посадову особу, яка керуватиме політикою та процедурами кадрової безпеки = \textbf{adminsecurity\_\allowbreak{}officer}\newline Визначена періодичність перегляду та оновлення поточної політики кадрової безпеки = \textbf{default\_\allowbreak{}deny\_\allowbreak{}ruleabac\_\allowbreak{}rule\_\allowbreak{}1}\newline Є події, які вимагають перегляду та оновлення поточної політики кадрової безпеки = \textbf{default\_\allowbreak{}deny\_\allowbreak{}ruleabac\_\allowbreak{}rule\_\allowbreak{}1}\newline Визначено періодичність перегляду та оновлення поточних процедур кадрової безпеки = \textbf{щорічно} \\ \hline
117 & Перевірка персоналу (PS-3) &  & PS-3 & Визначені умови, що вимагають повторної перевірки осіб = \textbf{}\newline Визначена частота повторної перевірки осіб, для яких це показано; PS-03a. проходять особи перевірку перед тим, як надати їм доступ до системи; PS-03b.[01] проходять особи повторну перевірку відповідно до умови, що вимагають повторної перевірки; PS-03b.[02] проводиться повторна перевірка у випадках, коли це зазначено, частота = \textbf{adminsecurity\_\allowbreak{}officer} \\ \hline
118 & Звільнення персоналу (PS-4) &  & PS-4 & Визначено період часу, протягом якого забороняється доступ до системи = \textbf{30}\newline Визначені теми інформаційної безпеки для обговорення під час проведення співбесід; PS-04a. при звільненні працівника доступ до системи вимикається протягом часового періоду; PS-04b. припиняють дію або анулюють будь-які автентифікатори та облікові дані після припинення трудових відносин з окремими особами; PS-04c. проводяться при звільненні окремих працівників співбесіди, які включають обговорення питань інформаційної безпеки; PS-04d. отримується після звільнення особи все майно, пов'язане з безпекою організаційної системи; PS-04e. зберігається доступ до організаційної інформації та систем, які раніше перебували під контролем особи, що звільняється, після припинення нею трудових відносин = \textbf{adminsecurity\_\allowbreak{}officer} \\ \hline
119 & Переведення персоналу (PS-5) &  & PS-5 & Визначені дії, які мають бути ініційовані після переведення або перепризначення = \textbf{loginlogoutfailed\_\allowbreak{}attempt}\newline Визначено період часу, протягом якого мають бути здійснені дії з переведення або перепризначення після переведення або перепризначення = \textbf{loginlogoutfailed\_\allowbreak{}attempt}\newline Визначено персонал або ролі, про які необхідно повідомляти, коли осіб призначають на інші посади або переводять на інші посади в організації = \textbf{adminsecurity\_\allowbreak{}officer}\newline Визначено період часу, протягом якого необхідно повідомляти визначений організацією персонал або ролі, коли осіб перепризначають або переводять на інші посади в межах організації; PS-05a. переглядаються та підтверджуються поточні потреби в логічних та фізичних дозволах на доступ до систем та об'єктів при перепризначенні або переведенні осіб на інші посади в організації; PS-05b. були ініційовані дії з переведення або перепризначення протягом періоду часу після формальної дії з переведення; PS-05c. змінюється авторизація доступу за необхідності, щоб відповідати будь-яким змінам в оперативних потребах у зв'язку з перепризначенням або переведенням; PS-05d. було повідомлено персонал або ролі протягом часового періоду = \textbf{adminsecurity\_\allowbreak{}officer} \\ \hline
\multicolumn{5}{|>{\raggedright\arraybackslash}p{14.5cm}|}{\textbf{14. ОЦІНКА РИЗИКІВ (RA) (RA)}} \\ \hline
120 & id-spe-ra-1 &  & RA-1 & Визначено персонал або ролі, на які поширюється політика оцінювання ризику = \textbf{adminsecurity\_\allowbreak{}officer}\newline Визначено персонал або ролі, на які поширюються процедури, що сприяють здійсненню політики оцінювання ризику = \textbf{adminsecurity\_\allowbreak{}officer}\newline Визначена посадова особа, відповідальна за управління політикою та процедурами оцінювання ризику = \textbf{adminsecurity\_\allowbreak{}officer}\newline Визначена періодичність перегляду та оновлення поточної політики оцінювання ризику = \textbf{default\_\allowbreak{}deny\_\allowbreak{}ruleabac\_\allowbreak{}rule\_\allowbreak{}1}\newline Є події, які вимагають перегляду та оновлення поточної політики оцінювання ризику = \textbf{default\_\allowbreak{}deny\_\allowbreak{}ruleabac\_\allowbreak{}rule\_\allowbreak{}1}\newline Визначено періодичність перегляду поточних процедур оцінювання ризику = \textbf{щорічно} \\ \hline
121 & ОЦІНЮВАННЯ РИЗИКУ (RA-3) &  & RA-3 & Визначено періодичність оцінювання ризику = \textbf{щорічно}\newline Визначено персонал або ролі, до яких мають бути доведені результати оцінювання ризику = \textbf{adminsecurity\_\allowbreak{}officer}\newline Визначена періодичність оновлення оцінювання ризику; перегляду результатів RA-03a.01 проводиться оцінювання ризику для виявлення загроз і вразливостей у системі; RA-03a.02 проводиться оцінювання ризику для визначення ймовірності та розміру шкоди від несанкціонованого доступу, використання, розкриття, порушення, модифікації або знищення системи; інформації, яку вона обробляє, зберігає або передає; а також будь-якої пов'язаної з нею інформації; RA-03a.03 проводиться оцінювання ризику для визначення ймовірності та впливу несприятливих наслідків для фізичних осіб, що виникають у зв'язку з обробкою інформації, яка ідентифікує особу; RA-03b. інтегровані результати оцінювання ризику та рішення з управління ризиками з точки зору організації та місії або бізнес-процесів з оцінкою ризиків на системному рівні; RA-03c. результати оцінювання ризику задокументовані в ЗНАЧЕННЯ ВИБРАНОГО ПАРАМЕТРА; RA-03d. переглядаються результати 03\_\allowbreak{}ODP[03] частота>; RA-03e. поширюються результати оцінювання ризику серед персоналу або ролей; RA-03f. оновлюється оцінювання ризику частота або коли відбуваються значні зміни в системі, середовищі її функціонування або інших умовах, які можуть вплинути на стан безпеки або приватності системи оцінювання ризику <RA- = \textbf{adminsecurity\_\allowbreak{}officer} \\ \hline
122 & ОЦІНЮВАННЯ ПОСТАЧАННЯ (RA-3(1)) &  & RA-3(1) &  \\ \hline
123 & СКАНУВАННЯ ВРАЗЛИВОСТЕЙ (RA-5) &  & RA-5 & Визначена періодичність перевірки систем та розміщених на них застосунків на наявність вразливостей = \textbf{щорічно}\newline Визначено час реагування на усунення законних вразливостей відповідно до організаційної оцінення ризику = \textbf{30}\newline Потрібно ділитися інформацією, отриманою в процесі сканування вразливостей та оцінок контролю, з персоналом або ролями, з якими потрібно ділитися; RA-05a.[01] здійснюється моніторинг систем та розміщених застосунків на наявність вразливостей частота та/або випадковість відповідно до визначеного організацією процесу, а також коли виявляються та повідомляються нові вразливості, що потенційно можуть вплинути на систему; RA-05a.[02] перевіряються системи та розміщені застосунки на наявність вразливостей частота та/або випадковим чином відповідно до визначеного організацією процесу, а також коли виявляються та повідомляються нові вразливості, що потенційно можуть вплинути на систему; RA-05b. застосовуються інструменти та методи моніторингу вразливостей для забезпечення сумісності між інструментами; RA-05b.01 застосовуються інструменти та методи моніторингу вразливостей для автоматизації частини процесу управління вразливостями, використовуючи стандарти для переліку платформ, недоліків програмного забезпечення та неправильних конфігурацій; RA-05b.02 застосовуються інструменти та методи моніторингу вразливостей для полегшення взаємодії між інструментами та автоматизації частини процесу управління вразливостями шляхом використання стандартів для формування контрольних списків та процедур тестування; RA-05b.03 застосовуються інструменти та методи моніторингу вразливостей для полегшення взаємодії між інструментами та автоматизації частин процесу управління вразливостями шляхом використання стандартів для вимірювання впливу вразливостей; RA-05c. аналізуються звіти про сканування вразливостей та результати моніторингу вразливостей; RA-05d. усуваються легітимні вразливості час реагування відповідно до організаційної оцінки ризиків; RA-05e. надається інформація, отримана в процесі моніторингу вразливостей та оцінки контролю, персонал або ролі, щоб допомогти усунути подібні вразливості в інших системах; RA-05f. використовуються інструменти моніторингу вразливостей, які передбачають можливість швидкого оновлення вразливостей, що підлягають скануванню = \textbf{adminsecurity\_\allowbreak{}officer} \\ \hline
124 & id-spe-ra-5-2 &  & RA-5(2) & Визначено час реагування на усунення законних вразливостей відповідно до організаційної оцінення ризику = \textbf{30} \\ \hline
125 & РЕАГУВАННЯ НА РИЗИК (RA-7) &  & RA-7 &  \\ \hline
\multicolumn{5}{|>{\raggedright\arraybackslash}p{14.5cm}|}{\textbf{15. ПРИДБАННЯ СИСТЕМ ТА ПОСЛУГ (SA) (SA)}} \\ \hline
126 & id-spe-sa-1 &  & SA-1 & Визначено персонал або ролі, на які поширюватиметься політика придбання систем і послуг = \textbf{adminsecurity\_\allowbreak{}officer}\newline Визначено персонал або ролі, на які поширюються процедури придбання систем і послуг = \textbf{adminsecurity\_\allowbreak{}officer}\newline Визначено посадову особу, яка керуватиме політикою та процедурами придбання систем і послуг = \textbf{adminsecurity\_\allowbreak{}officer}\newline Визначено періодичність перегляду та оновлення поточної політики придбання систем і послуг = \textbf{default\_\allowbreak{}deny\_\allowbreak{}ruleabac\_\allowbreak{}rule\_\allowbreak{}1}\newline Є події, які вимагають перегляду та оновлення поточної політики придбання систем і послуг = \textbf{default\_\allowbreak{}deny\_\allowbreak{}ruleabac\_\allowbreak{}rule\_\allowbreak{}1}\newline Визначено частоту, з якою переглядаються та оновлюються поточні процедури придбання систем і послуг = \textbf{30} \\ \hline
127 & id-spe-sa-8 &  & SA-8 &  \\ \hline
128 & ЗОВНІШНІ ПОСЛУГИ ДЛЯ СИСТЕМИ (SA-9) &  & SA-9 & Визначені засоби контролю, які будуть застосовуватися зовнішніми постачальниками системних послуг = \textbf{автоматизований засіб моніторингу}\newline Визначені процеси, методи та техніки, що застосовуються для моніторингу дотримання вимог контролю зовнішніми постачальниками послуг; SA-09a.[01] дотримуються постачальники зовнішніх системних послуг вимог організаційної безпеки; SA-09a.[02] дотримуються постачальники зовнішніх системних послуг вимог приватності організації; SA-09a.[03] використовують постачальники зовнішніх системних послуг елементи керування; SA-09b.[01] визначено та задокументовано організаційний нагляд за зовнішніми системними послугами; SA-09b.[02] визначені та задокументовані ролі та обов'язки користувачів щодо зовнішніх системних сервісів; SA-09c. визначені та задокументовані ролі та обов'язки користувачів щодо зовнішніх системних послуг = \textbf{adminsecurity\_\allowbreak{}officer} \\ \hline
129 & id-spe-sa-22 &  & SA-22 &  \\ \hline
\multicolumn{5}{|>{\raggedright\arraybackslash}p{14.5cm}|}{\textbf{16. ЗАХИСТ СИСТЕМ ТА КОМУНІКАЦІЙ (SC) (SC)}} \\ \hline
130 & id-spe-sc-1 &  & SC-1 & Визначено персонал або ролі, до яких має бути доведена політика захисту системи та комунікацій = \textbf{adminsecurity\_\allowbreak{}officer}\newline Визначено персонал або ролі, на які поширюються процедури захисту системи та комунікацій = \textbf{adminsecurity\_\allowbreak{}officer}\newline Визначено посадову особу, яка керуватиме політикою та процедурами захисту системи та комунікацій = \textbf{adminsecurity\_\allowbreak{}officer}\newline Визначена періодичність перегляду та оновлення поточної політики захисту системи та комунікацій = \textbf{default\_\allowbreak{}deny\_\allowbreak{}ruleabac\_\allowbreak{}rule\_\allowbreak{}1}\newline Є події, які вимагають перегляду та оновлення поточної політики захисту системи та комунікацій = \textbf{default\_\allowbreak{}deny\_\allowbreak{}ruleabac\_\allowbreak{}rule\_\allowbreak{}1}\newline Визначена періодичність перегляду та оновлення поточних процедур захисту системи та засобів зв'язку = \textbf{щорічно} \\ \hline
131 & id-spe-sc-4 &  & SC-4 &  \\ \hline
132 & ДОСТУПНІСТЬ РЕСУРСІВ (SC-7) &  & SC-7 &  \\ \hline
133 & id-spe-sc-7-5 &  & SC-7(5) &  \\ \hline
134 & id-spe-sc-8 &  & SC-8 &  \\ \hline
135 & id-spe-sc-8-1 &  & SC-8(1) &  \\ \hline
136 & ВІДКЛЮЧЕННЯ МЕРЕЖІ (SC-10) &  & SC-10 & Визначено період бездіяльності, після якого система розриває мережеве з'єднання, пов'язане з сеансом зв'язку; SC08(05)\_\allowbreak{}ODP[02] мережеве з'єднання, пов'язане з сеансом зв'язку, розірвано в кінці сеансу або після період часу бездіяльності = \textbf{30} \\ \hline
137 & ВСТАНОВЛЕННЯ КЛЮЧАМИ (SC-12) &  & SC-12 & Встановлюються криптографічні ключі, коли в системі використовується криптографія відповідно до < SC-12\_\allowbreak{}ODP вимог > = \textbf{AES-256-GCM}\newline Здійснюється управління криптографічними ключами, коли в системі використовується криптографія, відповідно до вимог  = \textbf{AES-256-GCM} \\ \hline
138 & КРИПТОГРАФІЧНИЙ ЗАХИСТ (SC-13) &  & SC-13 & КРИПТОГРАФІЧНИЙ ЗАХИСТ МЕТА ОЦІНКИ: Визначити, чи: = \textbf{AES-256-GCM}\newline Визначено використання криптографічних засобів = \textbf{AES-256-GCM}\newline Визначено типи криптографії для кожного вказаного криптографічного використання; SC-13a. ідентифіковано використання>; SC-13b. реалізовано типи криптографії для кожного вказаного криптографічного використання (визначеного в SC-13\_\allowbreak{}ODP[01]). криптографічне <SC-13\_\allowbreak{}ODP[01] = \textbf{AES-256-GCM} \\ \hline
139 & id-spe-sc-15 &  & SC-15 &  \\ \hline
140 & МОБІЛЬНИЙ КОД (SC-18) &  & SC-18 &  \\ \hline
141 & АВТЕНТИФІКАЦІЯ СЕСІЇ (SC-23) &  & SC-23 &  \\ \hline
142 & ЗАХИСТ ІНФОРМАЦІЇ В СТАНІ СПОКОЮ (SC-28) &  & SC-28 &  \\ \hline
143 & id-spe-sc-28-1 &  & SC-28(1) & Реалізовані криптографічні механізми для запобігання несанкціонованому розкриттю інформації, що знаходиться в стані спокою на системних компонентах або носіях = \textbf{AES-256-GCM}\newline Реалізовані криптографічні механізми для запобігання несанкціонованій модифікації інформації, що знаходиться в стані спокою на системних компонентах або носіях = \textbf{AES-256-GCM} \\ \hline
\multicolumn{5}{|>{\raggedright\arraybackslash}p{14.5cm}|}{\textbf{17. ЦІЛІСНІСТЬ СИСТЕМИ ТА ІНФОРМАЦІЇ (SI) (SI)}} \\ \hline
144 & id-spe-si-1 &  & SI-1 & Визначено персонал або ролі, до яких має бути доведена політика цілісності системи та інформації = \textbf{adminsecurity\_\allowbreak{}officer}\newline Визначено персонал або ролі, на які поширюються процедури цілісності системи та інформації  = \textbf{adminsecurity\_\allowbreak{}officer}\newline Визначено посадову особу, відповідальну за управління системою та політикою і процедурами цілісності інформації = \textbf{adminsecurity\_\allowbreak{}officer}\newline Визначено періодичність перегляду та оновлення поточної політики цілісності системи та інформації = \textbf{default\_\allowbreak{}deny\_\allowbreak{}ruleabac\_\allowbreak{}rule\_\allowbreak{}1}\newline Є події, які вимагають перегляду та оновлення поточної політики цілісності системи та інформації = \textbf{default\_\allowbreak{}deny\_\allowbreak{}ruleabac\_\allowbreak{}rule\_\allowbreak{}1}\newline Визначено частоту, з якою переглядаються та оновлюються поточні цілісності системи та інформації = \textbf{30} \\ \hline
145 & ВИПРАВЛЕННЯ ДЕФЕКТІВ (SI-2) &  & SI-2 & Визначено період часу, протягом якого необхідно встановити оновлення програмного забезпечення, пов'язані з безпекою, після виходу оновлень; SI-02a.[01] виявлено недоліки системи; SI-02a.[02] повідомляється про недоліки системи; SI-02a.[03] виправлені недоліки системи; SI-02b.[01] перевіряються оновлення програмного забезпечення, пов'язані з усуненням недоліків, на ефективність перед встановленням; SI-02b.[02] перевіряються оновлення програмного забезпечення, пов'язані з виправленням дефектів, на наявність потенційних побічних ефектів перед встановленням; SI-02b.[03] перевіряються оновлення прошивки, пов'язані з усуненням недоліків, на ефективність перед встановленням; SI-02b.[04] перевіряються оновлення прошивки, пов'язані з усуненням недоліків, на наявність потенційних побічних ефектів перед встановленням; SI-02c.[01] встановлено оновлення програмного забезпечення, що стосуються безпеки, протягом часовий проміжок з моменту випуску оновлень; SI-02c.[02] встановлено оновлення мікропрограми, що стосуються безпеки, протягом часового періоду з моменту випуску оновлень; SI-02d. включено відновлення порушених прав у процес управління організаційною конфігурацією = \textbf{30} \\ \hline
146 & ЗАХИСТ ВІД ШКІДЛИВОГО КОДУ (SI-3) &  & SI-3 & Визначено частоту, з якою механізми шкідливого коду виконують сканування = \textbf{30}\newline Визначено дії, яких слід вжити у відповідь на виявлення шкідливого коду (якщо вибрано) = \textbf{loginlogoutfailed\_\allowbreak{}attempt} \\ \hline
147 & МОНІТОРИНГ СИСТЕМИ (SI-4) &  & SI-4 & Визначені методи та способи, що використовуються для виявлення несанкціонованого використання системи = \textbf{adminsecurity\_\allowbreak{}officer}\newline Визначена інформація про моніторинг системи, яка повинна надаватися персоналу або ролям = \textbf{adminsecurity\_\allowbreak{}officer}\newline Визначено персонал або ролі, яким має надаватися інформація про моніторинг системи = \textbf{adminsecurity\_\allowbreak{}officer}\newline Вибрано одне або декілька з наступних ЗНАЧЕНЬ ПАРАМЕТРІВ: \{за потребою; частота\} = \textbf{щорічно} \\ \hline
148 & id-spe-si-4-4 &  & SI-4(4) & Визначені критерії незвичної або несанкціонованої діяльності або умови для вхідного трафіку зв'язку = \textbf{}\newline Визначені критерії незвичайної або несанкціонованої діяльності або умови для вихідного трафіку зв'язку = \textbf{}\newline Здійснюється моніторинг вхідного комунікаційного трафіку частота на предмет незвичних або несанкціонованих дій або умов = \textbf{щорічно}\newline Контролюється вихідний трафік зв'язку частота на предмет незвичних або несанкціонованих дій або умов. вихідного виявлення = \textbf{щорічно} \\ \hline
149 & id-spe-si-5 &  & SI-5 & Вибрано одне або декілька з наступних ЗНАЧЕНЬ ПАРАМЕТРІВ: \{персонал або ролі; елементи; зовнішні організації\} = \textbf{adminsecurity\_\allowbreak{}officer}\newline Визначено персонал або ролі, до яких мають бути доведені попередження, поради та директиви з безпеки (якщо визначено) = \textbf{adminsecurity\_\allowbreak{}officer} \\ \hline
150 & id-spe-si-12 &  & SI-12 & Здійснюється управління інформацією в системі відповідно до чинних законів, наказів, директив, положень, політик, стандартів, інструкцій та експлуатаційних вимог = \textbf{default\_\allowbreak{}deny\_\allowbreak{}ruleabac\_\allowbreak{}rule\_\allowbreak{}1}\newline Зберігається інформація в системі відповідно до чинних законів, указів Президента, директив, положень, політик, стандартів, інструкцій та експлуатаційних вимог = \textbf{default\_\allowbreak{}deny\_\allowbreak{}ruleabac\_\allowbreak{}rule\_\allowbreak{}1}\newline Управління інформацією, що виводиться з системи, здійснюється відповідно до чинних законів, указів Президента, директив, положень, політик, стандартів, інструкцій та операційних вимог = \textbf{default\_\allowbreak{}deny\_\allowbreak{}ruleabac\_\allowbreak{}rule\_\allowbreak{}1}\newline Зберігається інформація, що виводиться з системи, відповідно до чинних законів, наказів, директив, положень, політик, стандартів, інструкцій та експлуатаційних вимог = \textbf{default\_\allowbreak{}deny\_\allowbreak{}ruleabac\_\allowbreak{}rule\_\allowbreak{}1} \\ \hline
\multicolumn{5}{|>{\raggedright\arraybackslash}p{14.5cm}|}{\textbf{18. id-spe-sr (SR)}} \\ \hline
151 & id-spe-sr-1 &  & SR-1 & Визначено персонал або ролі, на які поширюється політика управління ризиками ланцюга постачання = \textbf{adminsecurity\_\allowbreak{}officer}\newline Визначено персонал або ролі, на які поширюються процедури управління ризиками ланцюга постачання = \textbf{adminsecurity\_\allowbreak{}officer}\newline Визначена посадова особа, відповідальна за розробку, документування та розповсюдження політики та процедур управління ризиками ланцюга постачання = \textbf{adminsecurity\_\allowbreak{}officer}\newline Визначена періодичність перегляду та оновлення поточної політики управління ризиками ланцюга постачання = \textbf{default\_\allowbreak{}deny\_\allowbreak{}ruleabac\_\allowbreak{}rule\_\allowbreak{}1}\newline Є події, які вимагають перегляду та оновлення поточної політики управління ризиками ланцюга постачання = \textbf{default\_\allowbreak{}deny\_\allowbreak{}ruleabac\_\allowbreak{}rule\_\allowbreak{}1}\newline Визначена періодичність перегляду та оновлення поточної процедури управління ризиками ланцюга постачання = \textbf{щорічно} \\ \hline
152 & id-spe-sr-2 &  & SR-2 &  \\ \hline
153 & id-spe-sr-3 &  & SR-3 & Визначено персонал ланцюга поставок, з яким необхідно координувати процес або процеси виявлення та усунення слабких місць або недоліків в елементах і процесах ланцюга постачання = \textbf{adminsecurity\_\allowbreak{}officer}\newline Визначені засоби контролю ланцюга постачання, що застосовуються для захисту від ризиків ланцюга постачання для системи, системного компонента або системної послуги, а також для обмеження шкоди або наслідків від подій, пов'язаних з ланцюгом постачання = \textbf{автоматизований засіб моніторингу}\newline Вибрано одне або декілька з наступних ЗНАЧЕНЬ ПАРАМЕТРІВ: \{плани безпеки та конфіденційності; план управління ризиками ланцюга постачання; документ\}; SR-03\_\allowbreak{}ODP[05] визначено документ, що ідентифікує обрані та впроваджені процеси та засоби контролю ланцюга постачання (якщо обрано); SR-03a.[01] запроваджено процес або процеси для виявлення та усунення слабких місць або недоліків в елементах та процесах ланцюга постачання; SR-03a.[02] процес або процеси виявлення та усунення слабких місць або недоліків в елементах та процесах ланцюга постачання системи або компонента системи координується/координуються з персоналом ланцюга постачання; SR-03b. застосовуються засоби контролю ланцюга постачання для захисту від ризиків ланцюга постачання для системи, системного компонента або системної послуги, а також для обмеження шкоди або наслідків від подій, пов'язаних з ланцюгом постачання; SR-03c. задокументовані обрані та впроваджені процеси та засоби контролю ланцюга постачання в ЗНАЧЕННЯ ВИБІРКОВОГО ПАРАМЕТРА(ів) = \textbf{adminsecurity\_\allowbreak{}officer}\newline Визначено документ, що ідентифікує обрані та впроваджені процеси та засоби контролю ланцюга постачання (якщо обрано); SR-03a.[01] запроваджено процес або процеси для виявлення та усунення слабких місць або недоліків в елементах та процесах ланцюга постачання; SR-03a.[02] процес або процеси виявлення та усунення слабких місць або недоліків в елементах та процесах ланцюга постачання системи або компонента системи координується/координуються з персоналом ланцюга постачання; SR-03b. застосовуються засоби контролю ланцюга постачання для захисту від ризиків ланцюга постачання для системи, системного компонента або системної послуги, а також для обмеження шкоди або наслідків від подій, пов'язаних з ланцюгом постачання; SR-03c. задокументовані обрані та впроваджені процеси та засоби контролю ланцюга постачання в ЗНАЧЕННЯ ВИБІРКОВОГО ПАРАМЕТРА(ів) = \textbf{adminsecurity\_\allowbreak{}officer} \\ \hline
154 & id-spe-sr-5 &  & SR-5 &  \\ \hline
155 & ОЦІНКА ПОСТАЧАЛЬНИКІВ (SR-6) &  & SR-6 & Оцінюються та аналізуються ризики, пов'язані з ланцюгом постачання, які стосуються постачальників або підрядників та систем, компонентів системи або системних послуг, які вони надають < SR-06\_\allowbreak{}ODP частота > = \textbf{щорічно}\newline Визначена періодичність оцінки та аналізу ризиків, пов'язаних з ланцюгом постачання, що стосуються постачальників або підрядників, а також систем, компонентів системи або системних послуг, які вони надають = \textbf{щорічно} \\ \hline


\end{longtable}

\end{document}
